\documentclass[a4paper,12pt,twoside]{memoir}

% Castellano
\usepackage[spanish,es-tabla]{babel}
\selectlanguage{spanish}
\usepackage[utf8]{inputenc}
\usepackage[T1]{fontenc}
\usepackage{lmodern} % scalable font
\usepackage{microtype}
\usepackage{placeins}
\usepackage{eurosym}

\RequirePackage{booktabs}
\RequirePackage[table]{xcolor}
\RequirePackage{xtab}
\RequirePackage{multirow}

% Links
\usepackage[colorlinks]{hyperref}
\hypersetup{
	colorlinks,
	linkcolor={blue!80!black},
	citecolor={blue!50!black},
	urlcolor={blue!80!black}
}

% Landscape
\usepackage{pdflscape}

% Multi-page tables using
\usepackage{longtable}
\usepackage{tabularx}

% Ecuaciones
\usepackage{amsmath}

% Rutas de fichero / paquete
\newcommand{\ruta}[1]{{\sffamily #1}}

% Párrafos
\nonzeroparskip


% Imagenes
\usepackage{graphicx}
\newcommand{\imagen}[2]{
	\begin{figure}[!h]
		\centering
		\includegraphics[width=0.9\textwidth]{#1}
		\caption{#2}\label{fig:#1}
	\end{figure}
	\FloatBarrier
}

\newcommand{\imagenAncho}[3]{
	\begin{figure}[H]
		\centering
		\includegraphics[width=#3\textwidth]{#1}
		\caption{#2}\label{fig:#1}
	\end{figure}
	\FloatBarrier
}

\newcommand{\imagenflotante}[2]{
	\begin{figure}%[!h]
		\centering
		\includegraphics[width=0.9\textwidth]{#1}
		\caption{#2}\label{fig:#1}
	\end{figure}
}



% El comando \figura nos permite insertar figuras comodamente, y utilizando
% siempre el mismo formato. Los parametros son:
% 1 -> Porcentaje del ancho de página que ocupará la figura (de 0 a 1)
% 2 --> Fichero de la imagen
% 3 --> Texto a pie de imagen
% 4 --> Etiqueta (label) para referencias
% 5 --> Opciones que queramos pasarle al \includegraphics
% 6 --> Opciones de posicionamiento a pasarle a \begin{figure}
\newcommand{\figuraConPosicion}[6]{%
  \setlength{\anchoFloat}{#1\textwidth}%
  \addtolength{\anchoFloat}{-4\fboxsep}%
  \setlength{\anchoFigura}{\anchoFloat}%
  \begin{figure}[#6]
    \begin{center}%
      \Ovalbox{%
        \begin{minipage}{\anchoFloat}%
          \begin{center}%
            \includegraphics[width=\anchoFigura,#5]{#2}%
            \caption{#3}%
            \label{#4}%
          \end{center}%
        \end{minipage}
      }%
    \end{center}%
  \end{figure}%
}

%
% Comando para incluir imágenes en formato apaisado (sin marco).
\newcommand{\figuraApaisadaSinMarco}[5]{%
  \begin{figure}%
    \begin{center}%
    \includegraphics[angle=90,height=#1\textheight,#5]{#2}%
    \caption{#3}%
    \label{#4}%
    \end{center}%
  \end{figure}%
}
% Para las tablas
\newcommand{\otoprule}{\midrule [\heavyrulewidth]}
%
% Nuevo comando para tablas pequeñas (menos de una página).
\newcommand{\tablaSmall}[5]{%
 \begin{table}[H]
  \begin{center}
   \rowcolors {2}{gray!35}{}
   \begin{tabular}{#2}
    \toprule
    #4
    \otoprule
    #5
    \bottomrule
   \end{tabular}
   \caption{#1}
   \label{tabla:#3}
  \end{center}
 \end{table}
}

%
%Para el float H de tablaSmallSinColores
\usepackage{float}

%
% Nuevo comando para tablas pequeñas (menos de una página).
\newcommand{\tablaSmallSinColores}[5]{%
 \begin{table}[H]
  \begin{center}
   \begin{tabularx}{\linewidth}{lXlX}
    \toprule
    #4
    \otoprule
    #5
    \bottomrule
   \end{tabularx}
   \caption{#1}
   \label{tabla:#3}
  \end{center}
 \end{table}
}

\newcommand{\tablaApaisadaSmall}[5]{%
\begin{landscape}
  \begin{table}
   \begin{center}
    \rowcolors {2}{gray!35}{}
    \begin{tabular}{#2}
     \toprule
     #4
     \otoprule
     #5
     \bottomrule
    \end{tabular}
    \caption{#1}
    \label{tabla:#3}
   \end{center}
  \end{table}
\end{landscape}
}

%
% Nuevo comando para tablas grandes con cabecera y filas alternas coloreadas en gris.
\newcommand{\tabla}[6]{%
  \begin{center}
    \tablefirsthead{
      \toprule
      #5
      \otoprule
    }
    \tablehead{
      \multicolumn{#3}{l}{\small\sl continúa desde la página anterior}\\
      \toprule
      #5
      \otoprule
    }
    \tabletail{
      \hline
      \multicolumn{#3}{r}{\small\sl continúa en la página siguiente}\\
    }
    \tablelasttail{
      \hline
    }
    \bottomcaption{#1}
    \rowcolors {2}{gray!35}{}
    \begin{xtabular}{#2}
      #6
      \bottomrule
    \end{xtabular}
    \label{tabla:#4}
  \end{center}
}

%
% Nuevo comando para tablas grandes con cabecera.
\newcommand{\tablaSinColores}[6]{%
  \begin{center}
    \tablefirsthead{
      \toprule
      #5
      \otoprule
    }
    \tablehead{
      \multicolumn{#3}{l}{\small\sl continúa desde la página anterior}\\
      \toprule
      #5
      \otoprule
    }
    \tabletail{
      \hline
      \multicolumn{#3}{r}{\small\sl continúa en la página siguiente}\\
    }
    \tablelasttail{
      \hline
    }
    \bottomcaption{#1}
    \begin{xtabular}{#2}
      #6
      \bottomrule
    \end{xtabular}
    \label{tabla:#4}
  \end{center}
}

%
% Nuevo comando para tablas grandes sin cabecera.
\newcommand{\tablaSinCabecera}[5]{%
  \begin{center}
    \tablefirsthead{
      \toprule
    }
    \tablehead{
      \multicolumn{#3}{l}{\small\sl continúa desde la página anterior}\\
      \hline
    }
    \tabletail{
      \hline
      \multicolumn{#3}{r}{\small\sl continúa en la página siguiente}\\
    }
    \tablelasttail{
      \hline
    }
    \bottomcaption{#1}
  \begin{xtabular}{#2}
    #5
   \bottomrule
  \end{xtabular}
  \label{tabla:#4}
  \end{center}
}



\definecolor{cgoLight}{HTML}{EEEEEE}
\definecolor{cgoExtralight}{HTML}{FFFFFF}

%
% Nuevo comando para tablas grandes sin cabecera.
\newcommand{\tablaSinCabeceraConBandas}[5]{%
  \begin{center}
    \tablefirsthead{
      \toprule
    }
    \tablehead{
      \multicolumn{#3}{l}{\small\sl continúa desde la página anterior}\\
      \hline
    }
    \tabletail{
      \hline
      \multicolumn{#3}{r}{\small\sl continúa en la página siguiente}\\
    }
    \tablelasttail{
      \hline
    }
    \bottomcaption{#1}
    \rowcolors[]{1}{cgoExtralight}{cgoLight}

  \begin{xtabular}{#2}
    #5
   \bottomrule
  \end{xtabular}
  \label{tabla:#4}
  \end{center}
}




\graphicspath{ {./img/} }

% Capítulos
\chapterstyle{bianchi}
\newcommand{\capitulo}[2]{
	\setcounter{chapter}{#1}
	\setcounter{section}{0}
	\chapter*{#2}
	\addcontentsline{toc}{chapter}{#2}
	\markboth{#2}{#2}
}

% Apéndices
\renewcommand{\appendixname}{Apéndice}
\renewcommand*\cftappendixname{\appendixname}

\newcommand{\apendice}[1]{
	%\renewcommand{\thechapter}{A}
	\chapter{#1}
}

\renewcommand*\cftappendixname{\appendixname\ }

% Formato de portada
\makeatletter
\usepackage{xcolor}
\newcommand{\tutor}[1]{\def\@tutor{#1}}
\newcommand{\course}[1]{\def\@course{#1}}
\definecolor{cpardoBox}{HTML}{E6E6FF}
\def\maketitle{
  \null
  \thispagestyle{empty}
  % Cabecera ----------------
\noindent\includegraphics[width=\textwidth]{cabecera}\vspace{1cm}%
  \vfill
  % Título proyecto y escudo informática ----------------
  \colorbox{cpardoBox}{%
    \begin{minipage}{.8\textwidth}
      \vspace{.5cm}\Large
      \begin{center}
      \textbf{TFG del Grado en Ingeniería Informática}\vspace{.6cm}\\
      \textbf{\LARGE\@title{}}
      \end{center}
      \vspace{.2cm}
    \end{minipage}

  }%
  \hfill\begin{minipage}{.20\textwidth}
    \includegraphics[width=\textwidth]{escudoInfor}
  \end{minipage}
  \vfill
  % Datos de alumno, curso y tutores ------------------
  \begin{center}%
  {%
    \noindent\LARGE
    Presentado por \@author{}\\ 
    en Universidad de Burgos --- \@date{}\\
    Tutor: \@tutor{}\\
  }%
  \end{center}%
  \null
  \cleardoublepage
  }
\makeatother

\newcommand{\nombre}{Lisa Cané Sáiz}
\newcommand{\nombreTFG}{Aplicación de Gestión de Aulas desarrollada con Vaadin y Spring Boot}

% Datos de portada
\title{\nombreTFG \\Documentación Técnica}
\author{\nombre}
\tutor{Jesús Manuel Maudes Raedo}
\date{\today}

\begin{document}

\maketitle


\cleardoublepage


%%%%%%%%%%%%%%%%%%%%%%%%%%%%%%%%%%%%%%%%%%%%%%%%%%%%%%%%%%%%%%%%%%%%%%%%%%%%%%%%%%%%%%%%



\frontmatter


\clearpage

% Indices
\tableofcontents

\clearpage

\listoffigures

\clearpage

\listoftables
\clearpage


\mainmatter

\appendix

\apendice{Plan de proyecto software}

\section{Introducción}
En este anexo se trata la planificación temporal llevada a cabo y también el estudio de viabilidad, en cuanto a parte económica y legal.

\section{Evolución temporal}
El desarrollo del proyecto se ha basado en la metodología ágil de Scrum en líneas generales, ya que se ha llevado a cabo por una única persona (en lugar de las 4 u 8 personas que normalmente trabajan) y se han realizado iteraciones (\textit{sprints}) de duración variable.

\subsection{Sprint 0 (17/10/2019 - 06/12/2019)}
Toma de contacto con el proyecto, planteamiento de las herramientas básicas con las que va a trabajar y búsqueda de diferentes alternativas.

Las tareas planteadas fueron:
\begin{itemize}
\tightlist
    \item Crear el repositorio de GitHub en el que alojar el proyecto.
    \item Elegir la herramienta de gestión del proyecto.
    \item Registrarse e investigar sobre Overleaf.
    \item Investigar las opciones de \textit{hosting} para el despliegue de la aplicación.
    \item Investigar Vaadin.
    \item Investigar Google Calendar API y buscar alternativas.
    \item Investigar Google Sign-In y buscar alternativas.
\end{itemize}

\subsection{Sprint 1 (06/12/2020 - 04/03/2020)}
Toma de contacto con la API de Google Calendar y la de Google Sign-In, y con la documentación en Overleaf.

Las tareas planteadas fueron:
\begin{itemize}
\tightlist
    \item Investigar la alternativa de Microsoft para la gestión del calendario y la autenticación.
    \item Crear la estructura de la documentación de \LaTeX en Overleaf a partir de la plantilla.
    \item Crear una aplicación de prueba para la API de Google Calendar.
\end{itemize}

\subsection{Sprint 2 (04/03/2020 - 06/04/2020)}
Últimas pruebas de toma de contacto con Google Calendar y Google Sign-In.

Las tareas planteadas fueron:
\begin{itemize}
\tightlist
    \item Crear una cuenta de Google para el proyecto.
    \item Crear una rama de GitHub para el proyecto de prueba.
    \item Crear un evento de Google Calendar.
    \item Controlar la duplicación de calendarios.
    \item Insertar el proyecto de prueba en la nueva rama.
\end{itemize}

\subsection{Sprint 3 (06/04/2020 - 15/04/2020)}
Planteamiento de los problemas surgidos en las pruebas de las APIs de Google y posibles soluciones, planteamiento de los requisitos de la aplicación y documentación de aspectos relevantes y problemas surgidos.

Las tareas planteadas fueron:
\begin{itemize}
\tightlist
    \item Migración de Vaadin 8 a Vaadin 10 para intentar solucionar el problema con Google Sign-In.
    \item Revisar las opciones de \textit{hosting} para el despliegue.
    \item Plantear los requisitos de la aplicación.
    \item Documentar los aspectos relevantes y problemas surgidos.
\end{itemize}

\subsection{Sprint 4 (15/04/2020 - 23/04/2020)}
Desarrollo de la documentación de especificación de requisitos de la aplicación (casos de uso, prototipo de ventanas, etc.), prueba de despliegue de la aplicación de prueba, investigación de las bases de datos permitidas con la opción de despliegue elegida y finalización de tareas. 

Las tareas planteadas fueron:
\begin{itemize}
\tightlist
    \item Resolver problemas con la integración de Google Sign-In.
    \item Probar el despliegue de la aplicación de prueba en GitHub Pages.
    \item Investigar las bases de datos permitidas con la opción de despliegue elegida.
    \item Documentar los aspectos relevantes y problemas surgidos.
    \item Crear la documentación de documentación de especificación de requisitos.
\end{itemize}

\subsection{Sprint 5 (23/04/2020 - 04/05/2020)}
Búsqueda del \textit{hosting} que permita subir la aplicación de prueba desarrollada con Vaadin de entre las planteadas y finalización de tareas.

Las tareas planteadas fueron:
\begin{itemize}
\tightlist
    \item Encontrar el hosting que permita subir la \textit{app} desarrollada con Vaadin (de entre las opciones planteadas).
    \item Documentar los aspectos relevantes y problemas surgidos.
\end{itemize}

\subsection{Sprint 6 (04/05/2020 - 11/05/2020)}
Revisión de problemas surgidos y búsqueda de soluciones, creación del proyecto Eclipse para desarrollar la aplicación final y finalización de tareas.

Las tareas planteadas fueron:
\begin{itemize}
\tightlist
    \item Buscar soluciones a los problemas con Google Sign-In.
    \item Resolver problemas del despliegue de la \textit{app} de prueba en Heroku.
    \item Crear el proyecto Eclipse para desarrollar la aplicación final. 
    \item Continuar con la búsqueda del hosting que permita subir la \textit{app} de Vaadin.
\end{itemize}

\subsection{Sprint 7 (11/05/2020 - 27/05/2020)}
Revisión de la especificación de requisitos y modificación de los mismos, y continuación de búsqueda de soluciones a los problemas con Google Sign-In. 

Las tareas planteadas fueron:
\begin{itemize}
\tightlist
    \item Modificar la documentación de especificación de requisitos.
    \item Continuar con la búsqueda del hosting que permita subir la \textit{app} de Vaadin y de la base de datos permitida.
    \item Configurar Heroku para establecer la conexión con PostgreSQL.
\end{itemize}

\subsection{Sprint 8 (27/05/2020 - 16/06/2020)}
Revisión de la especificación de requisitos y modificación de los mismos, creación del script de generación de la base de datos y desarrollo de la aplicación tras finalizar las pruebas. 

Las tareas planteadas fueron:
\begin{itemize}
\tightlist
    \item Modificar la documentación de especificación de requisitos.
    \item Crear el modelo Entidad-Relación.
    \item Crear el script SQL de generación de la base de datos.
    \item Crear la ventana Mantenimieno de Centros y Departamentos.
    \item Crear la ventana Histórico de Reservas.
    \item Crear la ventana Consulta de Reservas y Disponibilidad de Aulas.
    \item Crear la ventana Mantenimiento de Aulas.
\end{itemize}

\subsection{Sprint 9 (16/06/2020 - 25/06/2020)}
Modificación de la especificación de requisitos y del código tras separar la ventana de Consulta de Reserva y Disponibilidad de Aulas en dos, despliegue de la \textit{app} en Heroku, cambio de versión de Vaadin 15 a Vaadin 16, limpieza y organización de código, integración de algunas pruebas automáticas e implementación de ventanas.

Las tareas planteadas fueron:
\begin{itemize}
\tightlist
    \item Separar la ventana de Consulta de Reserva y Disponibilidad de Aulas en dos.
    \item Desplegar la \textit{app} en Heroku y generar su base de datos.
    \item Añadir filtro "Día de la semana" a las ventanas de consulta.
    \item Cambiar de versión de Vaadin 15 a Vaadin 16.
    \item Solucionar problemas en consulta de disponibilidad de aulas.
    \item Limpiar y organizar el código.
    \item Integrar pruebas automáticas para la consulta de disponibilidad de aulas.
    \item Crear la ventana de Reserva de Aulas.
    \item Crear la ventana de Gestión de Reservas.
\end{itemize}

\subsection{Sprint 10 (25/06/2020 - 03/07/2020)}
Continuación de la memoria en \LaTeX, incluyendo la especificación de requisitos y actualizando el modelo ER, y del desarrollo de las ventanas, modificación de la tabla \textit{historico\_reservas}, despliegue de la \textit{app} para prueba y nuevo intento de configuración de Google Calendar y Google Sign-In.

Las tareas planteadas fueron:
\begin{itemize}
\tightlist
    \item Modificar la documentación de especificación de requisitos.
    \item Crear el modelo Entidad-Relación.
    \item Crear el script SQL de generación de la base de datos.
    \item Crear la ventana Mantenimieno de Centros y Departamentos.
    \item Crear la ventana Histórico de Reservas.
    \item Crear la ventana Consulta de Reservas y Disponibilidad de Aulas.
    \item Crear la ventana Mantenimiento de Aulas.
    \item Integrar de Google Calendar y Google Sign-In.
    \item Modificar la tabla \textit{historico\_reservas} eliminando las FK.
    \item Actualizar el despliegue de la \textit{app} en Heroku con las nuevas ventanas.
\end{itemize}

\subsection{Sprint 11 (03/07/2020 - 09/07/2020)}
Descartada la integración de Google Calendar y Google Sign-In, implementación de la ventana de login con Spring Security, continuación con la mejora y arreglo de bugs en ventanas y con la documentación de \LaTeX.

Las tareas planteadas fueron:
\begin{itemize}
\tightlist
    \item Crear la ventana de Login y configurar la seguridad con Spring Security.
    \item Añadir la tabla \textit{usuario} y modificar el resto para incluir las FK.
    \item Modificar las tablas \textit{reserva} e \textit{historico\_reservas}.
    \item Crear la ventana Mantenimiento de Usuarios.
    \item Restringir la reserva y gestión de reservas.
\end{itemize}

\subsection{Sprint 12 (09/07/2020 - )}
Validación de las reservas y la gestión de reservas, implementación de la ventana perfil de usuario,continuación con la mejora y arreglo de bugs en ventanas y con la documentación de \LaTeX.

Las tareas planteadas fueron:
\begin{itemize}
\tightlist
    \item Validar las reservas que se realizan (creaciones y modificaciones).
    \item Mostrar solo las reservas de los centros y departamentos de los que es responsable el usuario logueado en Gestión de Reservas.
    \item Solucionar bugs en las ventanas.
    \item Actualizar el despliegue de la \textit{app}.
    \item Añadir el \textit{logger} y terminar el tratamiento de excepciones.
    \item Crear la ventana Perfil de Usuario.
    \item Actualizar la especificación de requisitos.
\end{itemize}

\section{Estudio de viabilidad}
\subsection{Viabilidad económica}
En este apartado se muestran los costes económicos que supondría desarrollar el proyecto de forma real.

\subsubsection{Coste del personal}
El proyecto ha sido desarrollado por una única persona durante 9 meses a tiempo parcial. Considerando un salario bruto mensual de 2000 \euro, el coste total del personal sería el siguiente:

\tablaSmallSinColores{Coste total del personal}{ l | c }{coste_personal}
{\textbf{Concepto} & \textbf{Coste} \\}{
    Salario mensual neto & 1266,67 \euro \\
    Retención del IRPF (11,69\%) & 167,63 \euro \\
    Seguridad Social (28,3\%) & 566 \euro \\
    Salario mensual bruto & 2000 \euro \\ \hline
    Coste total de los 9 meses & 18000 \euro \\
}

\subsubsection{Coste \textit{hardware}}
Los costes de esta sección hacen referencia al equipo utilizado en el desarrollo del proyecto, tratándose de un equipo valoreado en 750 \euro.

Teniendo en cuenta una amortización de 4 años, el coste de los 9 meses sería:
\tablaSmallSinColores{Coste total del \textit{hardware}}{ l | c  | c }{coste_hardware}
{\textbf{Concepto} & \textbf{Coste} & \textbf{Amortización}\\}{
    Ordenador portátil & 750 \euro & 140,625 \euro\\
    Coste total de los 9 meses & 750 \euro\\
}

\subsubsection{Coste \textit{software}}
Los costes de esta sección hacen referencia al software no gratuito utilizado en el proyecto, como es el sistema operativo de Windows y el coste del despliegue en Heroku, considerando el plan más básico de 21,87\euro (25 dólares).

Teniendo en cuenta una amortización de 2 años, el coste de los 9 meses sería:
\tablaSmallSinColores{Coste total del \textit{software}}{ l | c  | c }{coste_software}
{\textbf{Concepto} & \textbf{Coste} & \textbf{Amortización}\\}{
    Windows 10 & 145 \euro & 54,37 \euro \\ 
    Heroku  & 196,83 \euro & 73,81 \euro \\ \hline
    Coste total de los 9 meses & 341,83 \euro \\
}

\subsection{Viabilidad legal}
En esta sección se muestra un listado con las licencias de cada dependencia que se ha utilizado en el proyecto. Las dependencias del segundo grupo son las utilizadas específicamente para la parte de test.

\tablaSmallSinColores{Dependencias del proyecto}{ l | l | l }{dependencias}
{\textbf{Software} & \textbf{Licencia} \\}{
    Vaadin & Apache License 2.0 \\
    Spring Boot Starter & Apache License 2.0 \\
    Spring Boot Devtools & Apache License 2.0 \\
    WebDriverManager & Apache License 2.0 \\
    Spring Boot Starter Data JPA & Apache License 2.0 \\
    PostgreSQL JDBC Driver & BSD 2-clause (o FreeBSD) \\ 
    Spring Security Web & Apache License 2.0 \\
    Spring Security Config & Apache License 2.0 \\
    Spring Boot Maven Plugin & Apache License 2.0 \\
    Vaadin Maven Plugin & Apache License 2.0 \\
    Flow Server Production Mode & Apache License 2.0 \\
    Maven Failsafe Plugin & Apache License 2.0 \\ \hline
    Spring Boot Starter Test & Apache License 2.0 \\
    H2 Database Engine & EPL 1.0 - MPL 2.0\\
}

Haciendo una comparativa de las limitaciones de cada licencia, y siguiendo las recomendaciones de GNU sobre las licencias~\cite{licencias_gnu}, el proyecto queda protegido bajo la licencia GPL-3, ya que todas las licencias utilizadas son compatibles con la misma.
\apendice{Especificación de Requisitos}

\section{Introducción}

En este anexo se recogen los objetivos generales y los requisitos del proyecto. Detallando los requisitos globales del proyecto y los requisitos funcionales y casos de uso de la aplicación establecidos en el desarrollo.

También se muestra el planteamiento inicial de la aplicación con los prototipos de las ventanas.

\section{Objetivos generales}
El objetivo general del proyecto es desarrollar una aplicación web que permita:
\begin{itemize}
	\item Reservar aulas de la universidad.
	\item Gestionar las reservas registradas (modificar y eliminar). 
	\item Consultar las reservas registradas en el sistema y la disponibilidad de las aulas.
	\item Consultar el histórico de reservas (creaciones, modificaciones, eliminaciones).
	\item Administrar los centros y departamentos registrados y sus responsables.
	\item Administrar las aulas de los centros.
	\item Administrar los usuarios.
	\item Administrar el propio perfil de usuario.
\end{itemize}

\section{Planteamiento inicial}
El planteamiento inicial de la aplicación, tras la primera revisión, constaba de 8 requisitos funcionales:
\begin{itemize}
	\item RF-1 Login.
	\item RF-2 Reservar aulas. 
	\item RF-3 Gestionar reservas.
	\item RF-4 Consultar reservas.
	\item RF-5 Consultar disponibilidad de aulas.
	\item RF-6 Histórico de reservas.
	\item RF-7 Administrar centros y departamentos.
	\item RF-8 Administrar aulas.
\end{itemize}

El mayor cambio respecto a la versión actual está en el login, pues en un principio se iba a tratar con \textit{Google Sign-In} y actualmente se ha implementado un login con \textit{Spring Security}, lo cual ha implicado añadir la ventana con el formulario de login y la administración de usuarios. 

Otro cambio destacable está en las ventanas de Consulta de Reservas y Consulta de Disponibilidad de Aulas, ya que en el planteamiento inicial estaban englobadas en una única ventana. 

\subsection{RF-2 Reservar aulas}
Prototipo de ventana:
\imagen{Prototipo_reserva_aulas_principal}{Prototipo de ventana Reserva de Aulas - Principal}
\imagen{Prototipo_reserva_aulas_reserva_rango}{Prototipo de ventana Reserva de Aulas - Reserva por rango de fechas}

\subsection{RF-3 Gestionar reservas}
Prototipo de ventana:
\imagen{Prototipo_gestion_reservas_principal}{Prototipo de ventana Gestión de Reservas - Principal}
\imagen{Prototipo_gestion_reservas_listado}{Prototipo de ventana Gestión de Reservas - Listado de reservas}
\imagen{Prototipo_gestion_reservas_editar}{Prototipo de ventana Gestión de Reservas - Editar una reserva}

\subsection{RF-4 Consultar reservas - RF-5 Consultar disponibilidad de aulas }
Prototipo de ventana:
\imagen{Prototipo_consulta_visitante}{Prototipo de ventana Consulta de Reservas y Aulas - Visitante}
\imagen{Prototipo_consulta_usuario_logueado}{Prototipo de ventana Consulta de Reservas y Aulas - Usuario logueado}
\imagen{Prototipo_consulta_reserva_visitante}{Prototipo de ventana Consulta de Reservas y Aulas - Listado consulta de reservas - Visitante}
\imagen{Prototipo_consulta_reserva_usuario_logueado}{Prototipo de ventana Consulta de Reservas y Aulas - Listado consulta de reservas - Usuario logueado}
\imagen{Prototipo_consulta_disponibilidad_aulas_visitante}{Prototipo de ventana Consulta de Reservas y Aulas - Listado consulta de disponibilidad de aulas - Visitante}
\imagen{Prototipo_consulta_disponibilidad_aulas_usuario_logueado}{Prototipo de ventana Consulta de Reservas y Aulas - Listado consulta de disponibilidad de aulas - Usuario logueado}

\subsection{RF-6 Histórico de reservas}
Prototipo de ventana:
\imagen{Prototipo_historico_reservas}{Prototipo de ventana Histórico de reservas}

\subsection{RF-7 Administrar centros y departamentos}
Prototipo de ventana:
\imagen{Prototipo_mantenimiento_propietarios_principal}{Prototipo de ventana Mantenimiento de Centros y Departamentos}
\imagen{Prototipo_mantenimiento_propietarios_editar}{Prototipo de ventana Mantenimiento de Centros y Departamentos - Editar}

\subsection{RF-8 Administrar aulas}
Prototipo de ventana:
\imagen{Prototipo_mantenimiento_aulas_principal}{Prototipo de ventana Mantenimiento de Aulas - Selección del centro o departamento}
\imagen{Prototipo_mantenimiento_aulas_listado_aulas}{Prototipo de ventana Mantenimiento de Centros y Departamentos - Listado de aulas del centro o departamento}
\imagen{Prototipo_mantenimiento_aulas_editar}{Prototipo de ventana Mantenimiento de Centros y Departamentos - Editar aula del centro o departamento}

\section{Catálogo de requisitos}

Los requisitos funcionales y no funcionales del proyecto son los siguientes:

\subsection{Requisitos funcionales}
\begin{itemize}
	\item \textbf{RF-1 Login}: el sistema debe restringir el acceso a ciertas funcionalidades.
	    \begin{itemize}
	        \item Los responsables de centros y departamentos sólo tendrán acceso a la reserva de aulas y gestión de reservas.
	        \item El administrador tendrá acceso a la visualización del histórico de reservas y a la administración de centros y departamentos y de aulas.
	        \item Cualquier usuario tendrá acceso a la consulta de reservas y la consulta de disponibilidad de aulas.
	    \end{itemize}
	    
	\item \textbf{RF-2 Reservar aulas}: el sistema debe permitir al usuario con rol de responsable reservar un aula
        de los centros o departamentos que tiene bajo responsabilidad.
        
	\item \textbf{RF-3 Gestionar reservas}: el sistema debe permitir al usuario con rol de responsable gestionar las reservas registradas a partir de la fecha actual (i.e., bajas y modificaciones), de los centros o departamentos que tiene bajo responsabilidad.
        
	\item \textbf{RF-4 Consultar reservas}: el sistema debe permitir a cualquier usuario consultar las reservas registradas.
	\item 
	    \textbf{RF-5 Consultar disponibilidad de aulas}: el sistema debe permitir a cualquier usuario consultar la disponibilidad de las aulas.
	    
	\item \textbf{RF-6 Consultar histórico de reservas}: el sistema debe permitir al administrador consultar el histórico de reservas, en el que se muestran las creaciones, modificaciones y eliminaciones de reservas.
	    
	\item \textbf{RF-7 Administrar centros y departamentos}: el sistema debe permitir al administrador dar de alta, modificar y dar de baja los centros y departamentos.
	    
	\item \textbf{RF-8 Administrar aulas}: el sistema debe permitir al administrador dar de alta, modificar y dar de baja las aulas de los centros y departamentos.
	    
	\item \textbf{RF-9 Administrar usuarios}: el sistema debe permitir al administrador dar de alta, modificar (bloquear y desbloquear) y dar de baja los usuarios.
	    
	\item \textbf{RF-10 Administrar perfil de usuario}: el sistema debe permitir al administrador y a los usuarios con rol de responsable modificar la información de su perfil de usuario y ver los centros y departamentos que tienen bajo su propiedad.

\end{itemize}

\subsection{Requisitos no funcionales}
\begin{itemize}
    \item \textbf{RNF-1 Usabilidad}: la aplicación debe resultar sencilla de utilizar y entender.
    
    \item \textbf{RNF-2 Seguridad}: la aplicación debe autenticar el acceso a las funcionalidades restringidas y administrar de forma adecuada los datos sensibles (como las contraseñas).
    
    \item \textbf{RNF-3 Mantenibilidad}: debe resultar sencillo realizar modificaciones o añadir nuevas funcionalidades a la aplicación.
    
    \item \textbf{RNF-4 Diseño web adaptativo}: el diseño de la aplicación se debe adaptar al dispositivo que se está utilizando para visitarla.
\end{itemize}

\section{Especificación de requisitos} 

\subsection{Funcionalidades del producto}
Las \underline{funcionalidades del software} son:

\tablaSmallSinColores{Funcionalidades del producto software}{ c | l }{funcionalidades_software}
{\textbf{ID} & \textbf{Requisito funcional} \\}{
    RF-1 & Login  \\
	RF-2 & Reservar aulas  \\
	RF-3 & Gestionar reservas  \\
	RF-4 & Consultar reservas  \\
	RF-5 & Consultar disponibilidad de aulas  \\
	RF-6 & Histórico de reservas  \\
	RF-7 & Administrar centros y departamentos  \\
	RF-8 & Administrar aulas  \\
	RF-9 & Administrar usuarios  \\
	RF-10 & Administrar perfil usuario  \\
}

\subsection{Clases y características de usuarios}
Los \underline{usuarios que van a utilizar el producto} son:

\tablaSmallSinColores{Clases y características de usuarios}{ l | l }{clases_caracteristicas_usuarios}
{\textbf{Usuario} & \textbf{Funcionalidad} \\}{
    \multirow{2}{*}{Administrador} 
        & Principalmente administrar centros, departamentos, aulas y usuarios, y consultar el histórico de reservas \\
        & También puede consultar las reservas y la disponibilidad de aulas y administrar su perfil de usuario \\ \hline
        
     \multirow{2}{*}{Responsable} 
        & Principalmente reservar aulas y gestionar las reservas de su centro o departamento \\
        & También puede consultar las reservas y la disponibilidad de aulas y administrar su perfil de usuario \\ \hline
        
    Visitante & Consultar las las reservas y la disponibilidad de aulas \\
}

\begin{landscape}
\subsection{Diagrama de casos de uso}
\imagenAncho{Diagrama_Casos_Uso}{Diagrama de casos de uso}{1}
\end{landscape}

\subsection{Casos de uso}

\subsubsection{RF-1 Login}

\textbf{Descripción}: el sistema debe restringir el acceso a las ventanas de administración, al histórico de reservas, a la reserva de aulas y a la gestión de reservas a usuarios registrados (a través del login) y permitir a cualquiera (usuarios y visitantes) acceder a las ventanas de consulta.

\textbf{Prioridad}: alta.

\textbf{Pre-condiciones}: ninguna.

\textbf{Entradas}:
    \begin{itemize}
    \tightlist
        \item Correo del usuario.
        \item Contraseña.
    \end{itemize}

\textbf{Flujo básico}:
    \begin{itemize}
        \item Se introduce el correo y la contraseña del usuario y se pulsa en el botón ``\textbf{\textit{Log in}}'' para acceder a las ventanas restringidas.
        
            \begin{itemize}
            \tightlist
                \item El sistema comprobará que el correo y la contraseña introducidos son correctos, vaciando las entradas de texto en caso contrario.
                
                \item Si se pulsa en ``\textbf{\textit{He olvidado la contraseña}}'' se mostrará una notificación ``\textit{Póngase en contacto con el administrador (\textbf{correo del administrador})}''.
                
                \item Una vez se ha entrado, el usuario tiene acceso a las siguientes acciones según su perfil:
                
                \tablaSmallSinColores{Tipos de acceso al sistema según el usuario logueado}{ l | l }{tipos_acceso_usuario_logueado}
                    {\textbf{Usuario} & \textbf{Acceso} \\}{
                        \multirow{6}{*}{Administrador} 
                            & Administrar los centros y departamentos \\
                            & Administrar las aulas \\ 
                            & Administrar los usuarios \\ 
                            & Consultar el histórico de reservas \\
                            & Consultar las reservas y la disponibilidad de las aulas \\
                            & Administrar su perfil \\ \hline
                            
                         \multirow{4}{*}{Responsable} 
                            & Reservar aulas de su centro o departamento \\
                            & Gestionar las reservas (modificar y eliminar) de su centro o departamento \\ 
                            & Consultar las reservas y la disponibilidad de las aulas \\ 
                            & Administrar su perfil \\ \hline
                            
                        Visitante & No puede loguearse \\
                    }
                    
                \item El usuario puede cerrar su sesión haciendo clic en el enlace ``\textbf{\textit{Cerrar sesión}}'' que se encuentra en la parte superior derecha de la ventana.
            \end{itemize}
            
        \item Se pulsa en el link ``\textbf{\textit{Acceso consulta de reservas y disponibilidad de aulas}}'' para acceder a las ventanas de consulta sin Login.
    \end{itemize}

\textbf{Salida}: acceso a la ventana con el menú de opciones personalizado en función del usuario.

\textbf{Campos de la ventana}:
\tablaSmallSinColores{Campos de la ventana Login}{ l | l | l | l }{campos_ventana_login}
    {\textbf{Campo} & \textbf{Tipo} & \textbf{Obligatorio} & \textbf{Descripción}\\}{
        \textbf{Usuario} & Texto & Sí & Campo en el que se introduce el correo del usuario \\ \hline
            
        \textbf{Contraseña} & Texto & Sí & Campo en el que se introduce la contraseña del usuario \\
    }
    
\subsubsection{RF-2 Reservar Aulas}

\textbf{Descripción}: el sistema debe permitir al responsable reservar un aula de los centros o departamentos que tiene bajo responsabilidad.

\textbf{Prioridad}: alta.

\textbf{Pre-condiciones}: haber accedido al sistema como responsable.

\textbf{Entradas}:
    \begin{itemize}
    \tightlist
        \item Fecha inicio.
        \item Fecha fin*.
        \item Hora inicio.
        \item Hora fin.
        \item Centro o departamento.
        \item Aula.
        \item Día de la semana*.
        \item Motivo de la reserva.
        \item Persona a cargo de la reserva (A cargo de).
    \end{itemize}
    
* En caso de reservas por rango de fechas.

\textbf{Flujo básico}:
    \begin{enumerate}
        \item Se selecciona la opción ``\textbf{\textit{Reserva de Aulas}}'' del menú lateral.
        
        \item Si el responsable quiere hacer una reserva por rango de fechas, marca el check ``\textbf{\textit{Reserva por rango de fechas}}''.
            \begin{itemize}
                \item Cuando se marca, el sistema habilita los campos ``\textit{Fecha fin}'' y ``\textit{Día de la semana}'' y cambia el nombre del campo ``\textit{Fecha}'' a ``\textit{Fecha inicio}'', y viceversa.
            \end{itemize}
            
        \item Se rellenan todos los campos de la ventana y se hace clic en el botón ``\textbf{\textit{Reservar}}'.
            \begin{itemize}
            \tightlist
                \item El sistema comprobará que no se ha dejado ningún campo en blanco, mostrando un mensaje de alerta ``\textit{Todos los campos son obligatorios}'' en caso contrario.
                
                \item El sistema comprobará que la ``\textit{Fecha inicio}'' es menor que la ``\textit{Fecha fin}'' (en una reserva por rango de fechas), mostrando una notificación de error ``\textit{La fecha de inicio de la reserva debe ser menor que la fecha de fin}'' en caso contrario.
                
                \item El sistema comprobará que el aula está disponible en la/s fecha/s y horas indicadas, mostrando una notificación de error ``\textit{El aula no está disponible en la/s fecha/s y horas indicadas}''. 
            \end{itemize}
    \end{enumerate}

\textbf{Salida}: notificación de error en caso de haber introducido una fecha de inicio mayor que la fecha de fin (en reservas por rango de fechas), de haber introducido una hora de inicio mayor que la hora de fin o de que el aula no esté disponible.

\textbf{Campos de la ventana}:
\tablaSmallSinColores{Campos de la ventana Reserva de Aulas}{llll}{campos_ventana_reserva_aulas}
    {\textbf{Campo} & \textbf{Tipo} & \textbf{Obligatorio} & \textbf{Descripción}\\}{
        \textbf{Fecha inicio} & Fecha & Sí & Campo de selección única \\ \hline
        \textbf{Fecha fin} & Fecha & Sí * & Campo de selección única \\ \hline
        \textbf{Hora inicio} & Hora & Sí & Campo de texto / selección \\ \hline
        \textbf{Hora fin} & Hora & Sí & Campo de texto / selección \\ \hline
        \textbf{Centro/ Departamento} & Texto & Sí & Campo de selección única \\ \hline
        \textbf{Aula} & Texto & Sí & Campo de selección única \\ \hline
        \textbf{Día de la semana} & Texto & No & Campo de selección única \\ \hline
        \textbf{Motivo} & Texto & Sí & Campo de texto \\ \hline
        \textbf{A cargo de} & Texto & Sí & Campo de texto \\
    }
    
* En caso de reservas por rango de fechas.

\subsubsection{RF-3 Gestionar reservas}

\textbf{Descripción}: el sistema debe permitir al responsable gestionar las reservas registradas a partir de la fecha actual, de los centros o departamentos que tiene bajo responsabilidad.

\textbf{Prioridad}: alta.

\textbf{Pre-condiciones}: haber accedido al sistema como responsable.

\textbf{Entradas}:
    \begin{itemize}
    \tightlist
        \item Fecha desde*.
        \item Fecha hasta*.
        \item Hora desde*.
        \item Hora hasta*.
        \item Día de la semana*.
        \item Fecha.
        \item Hora inicio.
        \item Hora fin.
        \item Centro o departamento.
        \item Aula.
        \item Motivo de la reserva.
        \item Persona a cargo de la reserva (A cargo de).
    \end{itemize}
    
* Entradas del formulario de búsqueda de reservas.

\textbf{Flujo básico}:
    \begin{enumerate}
        \item Se selecciona la opción ``\textbf{\textit{Gestión de Reservas}}'' del menú lateral.
        
        \item El sistema muestra un listado con las reservas registradas a partir de la fecha actual de todos los centros y departamentos bajo responsabilidad del usuario logueado en una tabla.
        
            \begin{itemize}
            \tightlist
                \item Si no se ha realizado ninguna reserva a partir de la fecha actual, se mostrará una notificación ``\textit{No hay reservas realizadas a partir de la fecha actual o que cumplan con los filtros aplicados}''.
                
                \item  Para cada elemento de la tabla se muestra: fecha, día de la semana, hora de inicio, hora de fin, aula, centro en que se encuentra el aula, motivo de la reserva, persona a cargo de la reserva y el centro o departamento en nombre del cual se hace la reserva.
            \end{itemize}
        
        \item Se introducen los filtros que se quieran aplicar en los campos correspondientes y se hace clic en el botón ``\textbf{\textit{Buscar}}''.
        
            \begin{itemize}
            \tightlist
                \item El sistema comprobará que la ``\textit{Fecha desde}'' es menor que la ``\textit{Fecha hasta}'', mostrando una notificación de error ``\textit{La fecha desde la que se quiere filtrar debe ser menor que la fecha hasta la que se quiere filtrar}'' en caso contrario.
                
                \item El sistema comprobará que la ``\textit{Hora desde}'' es menor que la ``\textit{Hora hasta}'', mostrando una notificación de error ``\textit{La hora desde la que se quiere filtrar debe ser menor que la hora hasta la que se quiere filtrar}'' en caso contrario.
                
                \item Si no hay reservas que cumplan con los filtros aplicados, se mostrará una notificación ``\textit{No hay reservas realizadas a partir de la fecha actual o que cumplan con los filtros aplicados}''.
            \end{itemize}
            
        \item El responsable puede \textbf{limpiar los filtros aplicados} clicando en el botón ``\textbf{\textit{Limpiar filtros}}''. El sistema vaciará todos los campos de filtrado y actualizará el listado de reservas con aquellas registradas a partir de la fecha actual.
        
        \item El responsable puede \textbf{modificar una reserva} seleccionando la fila correspondiente y clicando en el botón ``\textbf{\textit{Modificar}}''.
        
            \begin{itemize}
            \tightlist
                \item El sistema comprobará que se ha seleccionado una única reserva, mostrando una notificación de error ``\textit{Solo se puede modificar una reserva}'' en caso de que se haya seleccionado más de una o ``\textit{No se ha seleccionado ninguna reserva}'' si no se ha seleccionado ninguna.
                
                \item El sistema muestra los datos actuales de la reserva en los campos correspondientes.
                
                \item El responsable modifica los campos.
                
                \item Puede guardar los cambios clicando en el botón ``\textbf{\textit{Guardar}}'' o descartarlos clicando en ``\textbf{\textit{Cancelar}}''.
                    \begin{itemize}
                    \tightlist
                        \item El sistema comprobará que no se ha dejado ningún campo en blanco, mostrando un mensaje de alerta ``\textit{Todos los campos son obligatorios}'' en caso contrario.
                        
                        \item El sistema comprobará que la ``\textit{Hora inicio}'' es menor que la ``\textit{Hora fin}'', mostrando una notificación de error ``\textit{La hora de inicio de la reserva debe ser menor que la hora de fin}'' en caso contrario.
                        
                        \item El sistema comprobará que el aula está disponible en la/s fecha/s y horas indicadas, mostrando una notificación de error ``\textit{El aula no está disponible en la/s fecha/s y horas indicadas}''.
                    \end{itemize}
            \end{itemize}
        
        \item El responsable puede \textbf{eliminar reservas} seleccionando las filas correspondientes, o marcando la casilla de la cabecera si quiere eliminar todas, y clicando en el botón ``\textbf{\textit{Eliminar}}''.
        
            \begin{itemize}
            \tightlist
                \item El sistema comprobará que se ha seleccionado al menos una reserva, mostrando una notificación de error ``\textit{No se ha seleccionado ninguna reserva}'' en caso contrario.
                
                \item El sistema mostrará un cuadro de diálogo con el mensaje ``\textit{¿Desea eliminar la/s reserva/s seleccionada/s definitivamente? Esta acción no se puede deshacer.}''.
                
                \item El responsable puede eliminarlas definitivamente clicando en ``\textbf{\textit{Confirmar}}'' o cancelar la acción clicando en ``\textbf{\textit{Cancelar}}''.
            \end{itemize}
            
        \item El responsable puede \textbf{ver la información acerca de la gestión de reservas} clicando en el botón con el icono de información situado a la derecha del botón ``\textbf{\textit{Eliminar}}''.
        
            \begin{itemize}
            \tightlist
                \item El sistema mostrará un cuadro de diálogo con el mensaje ``\textit{Sólo se pueden modificar las reservas de una en una. Si desea modificar una reserva de un rango de fechas debe eliminar todas las reservas correspondientes al rango y realizar una nueva reserva de rango.}''. Se puede cerrar clicando en el botón ``\textbf{\textit{Cerrar}}'' o en cualquier parte de la pantalla fuera del cuadro.
            \end{itemize}
    \end{enumerate}

\textbf{Salidas}: 
\begin{itemize}
\tightlist
    \item Listado con las reservas registradas que cumplan con los filtros aplicados, o por defecto con las reservas registradas por los centros o departamentos de los que es responsable el usuario logueado.
    
    \item Notificación de error en caso de haber dejado algún campo obligatorio en blanco, de haber introducido campos incorrectos, o de que el aula no esté disponible.
    
    \item Notificación en caso de que no haya reservas realizadas a partir de la fecha actual o que cumplan con los filtros aplicados.
    
    \item Cuadro de diálogo para confirmar la eliminación de la/s reserva/s.
    
    \item Información acerca de la gestión de reservas en caso de clicar en el botón de información.
\end{itemize}

\textbf{Campos de la ventana}:
\tablaSmallSinColores{Campos de la ventana Gestión de Reservas}{llll}{campos_ventana_gestion_reservas}
    {\textbf{Campo} & \textbf{Tipo} & \textbf{Obligatorio} & \textbf{Descripción}\\}{
        \textbf{Fecha desde*} & Fecha & No & Campo de selección única \\ \hline
        \textbf{Fecha hasta*} & Fecha & No * & Campo de selección única \\ \hline
        \textbf{Hora desde*} & Hora & No & Campo de texto / selección \\ \hline
        \textbf{Hora hasta*} & Hora & No & Campo de texto / selección \\ \hline
        \textbf{Día de la semana*} & Texto & No & Campo de selección única \\ \hline
        \textbf{Fecha} & Fecha & Sí & Campo de selección única \\ \hline
        \textbf{Hora inicio} & Hora & Sí & Campo de texto / selección \\ \hline
        \textbf{Hora fin} & Hora & Sí & Campo de texto / selección \\ \hline
        \textbf{Centro/ Departamento} & Texto & Sí & Campo de selección única \\ \hline
        \textbf{Aula} & Texto & Sí & Campo de selección única \\ \hline
        \textbf{Motivo} & Texto & Sí & Campo de texto \\ \hline
        \textbf{A cargo de} & Texto & Sí & Campo de texto \\
    }
    
* Campos del formulario de búsqueda de reservas.

\subsubsection{RF-4 Consultar reservas}

\textbf{Descripción}: el sistema debe permitir a cualquier usuario consultar las reservas registradas.

\textbf{Prioridad}: alta.

\textbf{Pre-condiciones}: ninguna.

\textbf{Entradas}:
    \begin{itemize}
    \tightlist
        \item Fecha desde.
        \item Fecha hasta.
        \item Hora desde.
        \item Hora hasta.
        \item Día de la semana.
        \item Centro o departamento que ha realizado la reserva.
    \end{itemize}

\textbf{Flujo básico}:
    \begin{enumerate}
        \item Se selecciona la opción ``\textit{Consulta de Reservas}'' del menú lateral.
        
        \item Se introducen los filtros que se quieran aplicar en los campos correspondientes y se hace clic en el botón ``\textbf{\textit{Buscar}}''.
            \begin{itemize}
            \tightlist
                \item El sistema comprobará que se ha seleccionado un centro o departamento, mostrando una notificación de error ``\textit{El Centro/Departamento es un campo obligatorio}'' en caso contrario.
                
                \item El sistema comprobará que la ``\textit{Fecha desde}'' es menor que la ``\textit{Fecha hasta}'', mostrando una notificación de error ``\textit{La fecha desde la que se quiere filtrar debe ser menor que la fecha hasta la que se quiere filtrar}'' en caso contrario.
                
                \item El sistema comprobará que la ``\textit{Hora desde}'' es menor que la ``\textit{Hora hasta}'', mostrando una notificación de error ``\textit{La hora desde la que se quiere filtrar debe ser menor que la hora hasta la que se quiere filtrar}'' en caso contrario.
            \end{itemize}
            
        \item El sistema muestra el listado con las reservas que cumplen con los filtros aplicados en una tabla.
            \begin{itemize}
            \tightlist
                \item Si no hay reservas que coincidan con los filtros aplicados, se mostrará una notificación ``\textit{No hay reservas realizadas a partir de la fecha actual o que cumplan con los filtros aplicados}''.
                
                \item Para cada elemento de la tabla se muestra: fecha, día de la semana, hora de inicio, hora de fin, aula, centro en que se encuentra el aula, motivo de la reserva, persona a cargo de la reserva y el usuario encargado de hacer la reserva seguido del centro o departamento en nombre del cual se hace la reserva.
            \end{itemize}
            
        \item El usuario puede \textbf{limpiar los filtros aplicados} clicando en el botón ``\textbf{\textit{Limpiar filtros}}''. El sistema vaciará todos los campos de filtrado y ocultará la tabla de reservas.
            
    \end{enumerate}

\textbf{Salidas}: 
 \begin{itemize}
\tightlist
    \item Listado con las reservas que cumplen con los filtros aplicados.
    
    \item Notificación de error en caso de haber dejado algún campo obligatorio vacío o de haber introducido campos incorrectos.
    
    \item Notificación en caso de que no haya reservas que cumplan con los filtros aplicados.
\end{itemize}

\textbf{Campos de la ventana}:
\tablaSmallSinColores{Campos de la ventana Consulta de Reservas}{llll}{campos_ventana_consulta_reserva}
    {\textbf{Campo} & \textbf{Tipo} & \textbf{Obligatorio} & \textbf{Descripción}\\}{
        \textbf{Fecha desde} & Fecha & No & Campo de selección única \\ \hline
        \textbf{Fecha hasta} & Fecha & No & Campo de selección única \\ \hline
        \textbf{Hora desde} & Hora & No & Campo de texto / selección \\ \hline
        \textbf{Hora hasta} & Hora & No & Campo de texto / selección \\ \hline
        \textbf{Día de la semana} & Texto & No & Campo de selección única \\ \hline
        \textbf{Centro/ Departamento} & Texto & Sí & Campo de selección única \\
    }

\subsubsection{RF-5 Consultar disponibilidad de aulas}

\textbf{Descripción}: el sistema debe permitir a cualquier usuario consultar la disponibilidad de las aulas de los centros y departamentos.

\textbf{Prioridad}: alta.

\textbf{Pre-condiciones}: ninguna.

\textbf{Entradas}:
    \begin{itemize}
    \tightlist
        \item Fecha desde.
        \item Fecha hasta.
        \item Hora desde.
        \item Hora hasta.
        \item Día de la semana.
        \item Capacidad del aula.
        \item Número de ordenadores del aula.
        \item Centro o departamento propietario del aula.
        \item Aula.
    \end{itemize}

\textbf{Flujo básico}:
    \begin{enumerate}
        \item Se selecciona la opción ``\textit{Consulta de Disponibilidad de Aulas}'' del menú lateral.
        
        \item Se introducen los filtros que se quieran aplicar en los campos correspondientes y se hace clic en el botón ``\textbf{\textit{Buscar}}''.
            \begin{itemize}
            \tightlist
                \item El sistema comprobará que se ha seleccionado un centro o departamento, mostrando una notificación de error ``\textit{El Centro/Departamento es un campo obligatorio}'' en caso contrario.
                
                \item El sistema comprobará que se han rellenado todos los filtros ``\textit{Fecha desde}'', ``\textit{Hora desde}'' y ``\textit{Hora hasta}'') si se ha introducido uno de ellos y/o ``\textit{Fecha hasta}'', mostrando una notificación de error ``\textit{Para filtrar por fecha y hora, ``Fecha desde'', ``Hora desde'' y ``Hora hasta'' son campos obligatorios}''. Si no se introduce la ``\textit{Fecha hasta}'', se toma el valor de ``\textit{Fecha desde}'' (para comprobaciones de un día concreto).
                
                \item El sistema comprobará que la ``\textit{Fecha desde}'' es menor que la ``\textit{Fecha hasta}'', mostrando una notificación de error ``\textit{La fecha desde la que se quiere filtrar debe ser menor que la fecha hasta la que se quiere filtrar}'' en caso contrario.
                
                \item El sistema comprobará que la ``\textit{Hora desde}'' es menor que la ``\textit{Hora hasta}'', mostrando una notificación de error ``\textit{La hora desde la que se quiere filtrar debe ser menor que la hora hasta la que se quiere filtrar}'' en caso contrario.
            \end{itemize}
            
        \item El sistema muestra el listado con las aulas disponibles que cumplen con los filtros aplicados en una tabla.
            \begin{itemize}
            \tightlist
                \item Si no hay aulas disponibles que coincidan con los filtros aplicados, se mostrará una notificación ``\textit{No hay aulas que cumplan con los filtros aplicados}''.
                
                \item Para cada elemento de la tabla se muestra: nombre del aula, centro en el que se encuentra el aula, capacidad y número de ordenadores.
            \end{itemize}
            
        \item El usuario puede \textbf{limpiar los filtros aplicados} clicando en el botón ``\textbf{\textit{Limpiar filtros}}''. El sistema vaciará todos los campos de filtrado y ocultará la tabla de aulas.
            
    \end{enumerate}

\textbf{Salidas}: 
 \begin{itemize}
\tightlist
    \item Listado con las aulas disponibles que cumplen con los filtros aplicados.
    
    \item Notificación de error en caso de haber dejado algún campo obligatorio vacío o de haber introducido campos incorrectos.
    
    \item Notificación en caso de que no haya aulas disponibles que cumplan con los filtros aplicados.
\end{itemize}

\textbf{Campos de la ventana}:
\tablaSmallSinColores{Campos de la ventana Consulta de Disponibilidad de aulas}{llll}{campos_ventana_consulta_disponibilidad_aulas}
    {\textbf{Campo} & \textbf{Tipo} & \textbf{Obligatorio} & \textbf{Descripción}\\}{
        \textbf{Fecha desde} & Fecha & No & Campo de selección única \\ \hline
        \textbf{Fecha hasta} & Fecha & No & Campo de selección única \\ \hline
        \textbf{Hora desde} & Hora & No & Campo de texto / selección \\ \hline
        \textbf{Hora hasta} & Hora & No & Campo de texto / selección \\ \hline
        \textbf{Día de la semana} & Texto & No & Campo de selección única \\ \hline
        \textbf{Capacidad} & Numérico & No & Campo de texto \\ \hline
        \textbf{Número de ordenadores} & Texto & No & Campo de texto \\ \hline
        \textbf{Centro/ Departamento} & Texto & Sí & Campo de selección única \\ \hline
        \textbf{Aula} & Texto & No & Campo de selección única \\ 
    }

\subsubsection{RF-6 Histórico de reservas}

\textbf{Descripción}: el sistema debe permitir al administrador consultar el histórico de reservas, en el que se muestran las creaciones, modificaciones y eliminaciones de reservas.

\textbf{Prioridad}: baja.

\textbf{Pre-condiciones}: haber accedido al sistema como administrador.

\textbf{Entradas}:
    \begin{itemize}
    \tightlist
        \item Fecha desde.
        \item Fecha hasta.
    \end{itemize}

\textbf{Flujo básico}:
    \begin{enumerate}
        \item Se selecciona la opción ``\textbf{\textit{Histórico de Reservas}}'' del menú lateral.
            
        \item El sistema muestra el listado con las operaciones realizadas sobre las reservas en una tabla. Por defecto se filtran las realizadas desde 7 días antes de la fecha actual hasta 7 días después.
            \begin{itemize}
            \tightlist
                \item Si en el rango de fechas no se ha realizado ninguna operación sobre las reservas se mostrará una notificación  ``\textit{No se han realizado operaciones con las reservas en esas fechas}''.
                
                \item Para cada elemento de la tabla se muestra: fecha de la operación, tipo de operación, motivo de la reserva, fecha de la reserva. hora de inicio de la reserva, hora de fin de la reserva y lugar de la reserva (aula - centro).
            \end{itemize}
            
        \item El administrador puede hacer \textbf{consultas por fecha} en el histórico introduciendo los filtros que se quieren aplicar y clicando en el botón ``\textbf{\textit{Buscar}}''.
            \begin{itemize}
            \tightlist
                \item El sistema comprobará que la ``\textit{Fecha desde}'' es menor que la ``\textit{Fecha hasta}'', mostrando una notificación de error ``\textit{La fecha desde la que se quiere filtrar debe ser menor que la fecha hasta la que se quiere filtrar}'' en caso contrario.
                
                \item Si en el rango de fechas no se ha realizado ninguna operación sobre las reservas se mostrará una notificación  ``\textit{No se han realizado operaciones con las reservas en esas fechas}''.
            \end{itemize}
            
        \item El administrador puede \textbf{limpiar los filtros aplicados} clicando en el botón ``\textbf{\textit{Limpiar filtros}}''. El sistema vaciará todos los campos de filtrado y ocultará la tabla de aulas.
    \end{enumerate}

\textbf{Salidas}: 
 \begin{itemize}
\tightlist
    \item Listado con las operaciones realizadas sobre las reservas.
    
    \item Notificación de error en caso de haber introducido campos incorrectos.
    
    \item Notificación en caso de que no haya operaciones realizadas sobre las reservas que cumplan con los filtros aplicados.
\end{itemize}

\textbf{Campos de la ventana}:
\tablaSmallSinColores{Campos de la ventana Histórico de Reservas}{llll}{campos_ventana_historico_reservas}
    {\textbf{Campo} & \textbf{Tipo} & \textbf{Obligatorio} & \textbf{Descripción}\\}{
        \textbf{Fecha desde} & Fecha & No & Campo de selección única \\ \hline
        \textbf{Fecha hasta} & Fecha & No & Campo de selección única \\ \hline
    }

\subsubsection{RF-7 Administrar centros y departamentos}

\textbf{Descripción}: el sistema debe permitir al administrador dar de alta, modificar y dar de baja los centros y departamentos.

\textbf{Prioridad}: baja.

\textbf{Pre-condiciones}: haber accedido como administrador.

\textbf{Entradas}:
    \begin{itemize}
    \tightlist
        \item ID del centro/departamento.
        \item Nombre del centro/departamento.
        \item Responsable.
    \end{itemize}

\textbf{Flujo básico}:
    \begin{enumerate}
        \item Se selecciona la opción ``\textbf{\textit{Mantenimiento de Centros y Departamentos}}'' del menú lateral.
        
        \item El sistema muestra el listado con los centros y departamentos registrados en una tabla.
            \begin{itemize}
                \item Para cada elemento de la tabla se muestra: nombre del centro o departamento, nombre y apellidos del responsable, correo del responsable y teléfono del responsable.
            \end{itemize}
            
        \item El administrador puede \textbf{añadir} centros y departamentos clicando el el botón ``\textbf{\textit{Añadir centro}}'' o ``\textbf{\textit{Añadir departamento}}'', respectivamente.
            \begin{itemize}
            \tightlist
                \item Se rellenan todos los campos de la ventana.
                
                \item Se registra el centro o departamento clicando en el botón ``\textbf{\textit{Guardar}}'' o se pueden descartar los cambios clicando en ``\textbf{\textit{Cancelar}}''. Si se ha dejado algún campo en blanco el botón `\textit{Guardar}'' permanecerá deshabilitado.
            \end{itemize}
            
        \item El administrador puede \textbf{modificar} los datos de un centro o departamento, o  \textbf{eliminarlo} clicando en la fila correspondiente.
            \begin{itemize}
            \tightlist
                \item El sistema muestra los datos actuales en los campos correspondientes.
                
                \item El administrador modifica los campos (excepto el ID que está deshabilitado).
                
            \item Puede guardar los cambios clicando en el botón `\textbf{\textit{Guardar}}'' o descartarlos clicando en ``\textbf{\textit{Cancelar}}''. Si se ha dejado algún campo en blanco el botón ``\textit{Guardar}'' permanecerá deshabilitado.
                
                \item Puede eliminar el centro o departamento clicando en el botón  ``\textbf{\textit{Eliminar}}''.
                    \begin{itemize}
                    \tightlist
                        \item El sistema comprobará si el centro o departamento que quiere eliminar tiene aulas asignadas, mostrando un cuadro de diálogo con un mensaje de confirmación ``\textit{Nombre del centro/departamento tiene aulas asignadas, ¿desea eliminarle definitivamente junto a todas sus aulas? Esta acción no se puede deshacer}'' en caso afirmativo o ``\textit{¿Desea eliminar nombre del centro/departamento definitivamente? Esta acción no se puede deshacer.}'' en caso contrario.
                        
                        \item El administrador puede eliminar definitivamente clicando en el botón  ``\textbf{\textit{Confirmar}}'' o cancelar la acción clicando en  ``\textbf{\textit{Cancelar}}''. Se mostrará una notificación ``\textit{Se ha eliminado nombre del centro/departamento correctamente}'' si se ha realizado con éxito.
                    \end{itemize}
            \end{itemize}
    \end{enumerate}

\textbf{Salidas}: 
\begin{itemize}
\tightlist
    \item Listado con los centros y departamentos registrados.
    
    \item Cuadro de diálogo para confirmar la eliminación de la/s reserva/s.
    
    \item Notificación de éxito si se ha eliminado correctamente.
\end{itemize}

\textbf{Campos de la ventana}:
\tablaSmallSinColores{Campos de la ventana Mantenimiento de Centros y Departamentos}{llll}{campos_ventana_mantenimiento_centros_departamentos}
    {\textbf{Campo} & \textbf{Tipo} & \textbf{Obligatorio} & \textbf{Descripción}\\}{
        \textbf{ID} & Texto & Sí & Campo de texto \\ \hline
        \textbf{Nombre} & Texto & Sí & Campo de texto \\ \hline
        \textbf{Responsable} & Texto & No & Campo de selección única \\ 
    }

\subsubsection{RF-8 Administrar aulas}

\textbf{Descripción}: el sistema debe permitir al administrador dar de alta, modificar y dar de baja las aulas de los centros y departamentos.

\textbf{Prioridad}: baja.

\textbf{Pre-condiciones}: haber accedido como administrador.

\textbf{Entradas}:
    \begin{itemize}
    \tightlist
        \item Centro o departamento (propietario) del que administrar las aulas.
        \item Nombre del aula.
        \item Capacidad.
        \item Número de ordenadores.
        \item Centro.
        \item Propietario del aula (centro o departamento).
    \end{itemize}

\textbf{Flujo básico}:
    \begin{enumerate}
        \item Se selecciona la opción ``\textbf{\textit{Mantenimiento de Aulas}}'' del menú lateral.
        
        \item Se selecciona el centro o departamento del que se quieren modificar y/o eliminar aulas.
        
        \item El sistema muestra el listado con las aulas del centro o departamento seleccionado en una tabla.
            \begin{itemize}
                \item Para cada elemento de la tabla se muestra: nombre del aula, centro en el que se encuentra, capacidad y número de ordenadores.
            \end{itemize}
            
        \item El administrador puede \textbf{añadir} aulas clicando el el botón ``\textbf{\textit{Añadir aula}}''.
            \begin{itemize}
            \tightlist
                \item Se rellenan todos los campos de la ventana.
                
                \item Se registra el aula clicando en el botón ``\textbf{\textit{Guardar}}'' o se pueden descartar los cambios clicando en  ``\textbf{\textit{Cancelar}}''. 
                    \begin{itemize}
                    \tightlist
                        \item Si se ha dejado algún campo en blanco el botón ``\textit{Guardar}'' permanecerá deshabilitado.
                        
                        \item El sistema comprobará que no existe un aula con el nombre introducido en el centro seleccionado, mostrando una notificación de error ``\textit{Ya existe un aula con ese nombre en el centro que ha seleccionado}'' en caso contrario.
                    \end{itemize}
            \end{itemize}
            
        \item El administrador puede \textbf{modificar} los datos de un aula, o  \textbf{eliminarla} clicando en la fila correspondiente.
            \begin{itemize}
            \tightlist
                \item El sistema muestra los datos actuales en los campos correspondientes.
                
                \item El administrador modifica los campos.
                
                \item Puede guardar los cambios clicando en el botón ``\textbf{\textit{Guardar}}'' o descartarlos clicando en  ``\textbf{\textit{Cancelar}}''. 
                    \begin{itemize}
                    \tightlist
                        \item Si se ha dejado algún campo en blanco el botón ``\textit{Guardar}'' permanecerá deshabilitado.
                        
                        \item El sistema comprobará que no existe un aula con el nombre introducido en el centro seleccionado, mostrando una notificación de error ``\textit{Ya existe un aula con ese nombre en el centro que ha seleccionado}'' en caso contrario.
                    \end{itemize}
                
                \item Puede eliminar el aula clicando en el botón  ``\textbf{\textit{Eliminar}}''.
                    \begin{itemize}
                    \tightlist
                        \item El sistema comprobará si el aula tiene reservas asociadas a partir de la fecha actual, mostrando un cuadro de diálogo con un mensaje de confirmación ``\textit{El aula nombre del aula de nombre del centro tiene reservas asociadas, ¿desea eliminarla definitivamente junto a todas sus reservas? Esta acción no se puede deshacer.}'' en caso afirmativo o ``\textit{¿Desea eliminar nombre del aula de nombre del centro definitivamente? Esta acción no se puede deshacer.}'' en caso contrario.
                        
                        \item El administrador puede eliminar definitivamente clicando en el botón  ``\textbf{\textit{Confirmar}}'' o cancelar la acción clicando en  ``\textbf{\textit{Cancelar}}''. Se mostrará un mensaje ``\textit{Se ha eliminado el aula nombre del aula correctamente}'' si se ha realizado con éxito.
                    \end{itemize}
            \end{itemize}
    \end{enumerate}

\textbf{Salidas}: 
\begin{itemize}
\tightlist
    \item Listado con las aulas del centro o departamento seleccionado.
    
    \item Cuadro de diálogo para confirmar la eliminación del aula.
    
    \item Mensaje de éxito si se ha eliminado correctamente.
\end{itemize}

\textbf{Campos de la ventana}:
\tablaSmallSinColores{Campos de la ventana Mantenimiento de Aulas}{llll}{campos_ventana_mantenimiento_aulas}
    {\textbf{Campo} & \textbf{Tipo} & \textbf{Obligatorio} & \textbf{Descripción}\\}{
        \textbf{Centro/Departamento} & Texto & No & Campo de selección única \\ \hline
        \textbf{Nombre del aula} & Texto & Sí & Campo de texto \\ \hline
        \textbf{Capacidad} & Numérico & Sí & Campo de texto \\ \hline
        \textbf{Número de ordenadores} & Numérico & Sí & Campo de texto \\ \hline
        \textbf{Centro} & Texto & Sí & Campo de selección única \\ \hline
        \textbf{Propietario del aula} & Texto & Sí & Campo de selección única \\ 
    }

\subsubsection{RF-9 Administrar usuarios}

\textbf{Descripción}: el sistema debe permitir al administrador dar de alta, modificar (bloquear y desbloquear) y dar de baja los usuarios.

\textbf{Prioridad}: baja.

\textbf{Pre-condiciones}: haber accedido como administrador.

\textbf{Entradas}:
    \begin{itemize}
    \tightlist
        \item Nombre del usuario.
        \item Apellidos del usuario.
        \item Correo del usuario.
        \item Contraseña.
        \item Teléfono del usuario.
        \item Rol del usuario.
        \item Si está o no bloqueado.
    \end{itemize}

\textbf{Flujo básico}:
    \begin{enumerate}
        \item Se selecciona la opción ``\textbf{\textit{Mantenimiento de Usuarios}}'' del menú lateral.
        
        \item El sistema muestra el listado con los usuarios registrados en una tabla.
            \begin{itemize}
                \item Para cada elemento de la tabla se muestra: nombre y apellidos del usuario, correo, teléfono y rol.
            \end{itemize}
            
        \item El administrador puede \textbf{añadir} usuarios clicando el el botón ``\textbf{\textit{Añadir usuario}}''.
            \begin{itemize}
            \tightlist
                \item Se rellenan todos los campos de la ventana.
                
                \item Se registra el usuario clicando en el botón ``\textbf{\textit{Guardar}}'' o se pueden descartar los cambios clicando en  ``\textbf{\textit{Cancelar}}''. 
                    \begin{itemize}
                    \tightlist
                        \item El sistema comprobará que el correo tiene el formato correcto, mostrando un mensaje de error bajo el campo en caso contrario.
                        
                        \item El sistema comprobará que la contraseña tiene como mínimo 5 caracteres, mostrando un mensaje de error ``\textit{Debe tener 5 o más caracteres}'' bajo el campo en caso contrario. La contraseña se almacena en la base de datos codificada (\textit{hasheada}).
                        
                        \item El sistema comprobará que el teléfono contiene 9 dígitos, mostrando un mensaje de error ``\textit{Debe tener 9 dígitos}'' bajo el campo en caso contrario.
                    \end{itemize}
            \end{itemize}
            
        \item El administrador puede \textbf{modificar} si el usuario está o no bloqueado, o  \textbf{eliminarlo} clicando en la fila correspondiente.
            \begin{itemize}
            \tightlist
                \item El sistema muestra los datos actuales en los campos correspondientes, manteniéndolos deshabilitados excepto el check de ``\textbf{\textit{Bloqueado}}''.
                
                \item El administrador bloquea o desbloquea al usuario.
                    \begin{itemize}
                    \tightlist
                        \item El sistema muestra un cuadro de diálogo con el mensaje ``\textit{¿Seguro que quiere desbloquear al usuario nombre y apellido del usuario?}'' en caso de que esté bloqueado, y viceversa.
                        
                        \item El administrador puede bloquear/desbloquear definitivamente pulsando en ``\textbf{\textit{Confirmar}}'' o cancelar la acción clicando en  ``\textbf{\textit{Cancelar}}''.
                    \end{itemize}
                
                
                \item Puede guardar los cambios clicando en el botón ``\textbf{\textit{Guardar}'}' o descartarlos clicando en  ``\textbf{\textit{Cancelar}}''. 
                
                \item Puede eliminar el usuario clicando en el botón  ``\textbf{\textit{Eliminar}}''.
                    \begin{itemize}
                    \tightlist
                        \item El sistema comprobará si el usuario que está eliminando es él mismo, el usuario logueado, mostrando una notificación de error ``\textit{No puede eliminar su propio usuario}'' en ese caso.
                        
                        \item El sistema comprobará si el usuario que está eliminando está bloqueado, mostrando una notificación de error ``\textit{El usuario está bloqueado y no se puede eliminar}'' en ese caso.
                        
                        \item El sistema comprobará si el usuario tiene centros o departamentos bajo su responsabilidad, mostrando una notificación de error ``\textit{El usuario nombre y apellidos del usuario tiene bajo su responsabilidad algún centro o departamento, reasígneles otro responsable primero}'' o ``\textit{¿Desea eliminar el usuario nombre y apellidos del usuario definitivamente? Esta acción no se puede deshacer}'' en caso contrario.
                        
                        \item Se puede eliminar definitivamente clicando en el botón  ``\textbf{\textit{Confirmar}}'' o cancelar la acción clicando en  ``\textbf{\textit{Cancelar}}''. Se mostrará un mensaje ``\textit{Se ha eliminado el usuario nombre y apellidos del usuario correctamente}'' si se ha realizado con éxito.
                    \end{itemize}
            \end{itemize}
    \end{enumerate}

\textbf{Salidas}: 
\begin{itemize}
\tightlist
    \item Listado con los usuarios registrados.
    
    \item Notificación de error en caso de intentar eliminar el usuario logueado o un usuario bloqueado.
    
    \item Cuadro de diálogo para confirmar la eliminación del aula.
    
    \item Mensaje de éxito si se ha eliminado correctamente.
\end{itemize}

\textbf{Campos de la ventana}:
\tablaSmallSinColores{Campos de la ventana Mantenimiento de Usuarios}{llll}{campos_ventana_mantenimiento_usuarios}
    {\textbf{Campo} & \textbf{Tipo} & \textbf{Obligatorio} & \textbf{Descripción}\\}{
        \textbf{Nombre} & Texto & Sí & Campo de texto \\ \hline
        \textbf{Apellidos} & Texto & Sí & Campo de texto \\ \hline
        \textbf{Correo} & Texto & Sí & Campo de e-mail \\ \hline
        \textbf{Contraseña} & Texto & Sí & Campo de texto \\ \hline
        \textbf{Teléfono} & Numérico & Sí & Campo numérico \\ \hline
        \textbf{Rol} & Texto & Sí & Campo de selección única \\ \hline
        \textbf{Bloqueado} & Booleano & No & Check \\
    }

\subsubsection{RF-10 Administrar perfil usuario}

\textbf{Descripción}:  el sistema debe permitir al administrador y a los responsables modificar la información de su perfil de usuario y ver los centros y departamentos que tienen bajo su propiedad.

\textbf{Prioridad}: baja.

\textbf{Pre-condiciones}: haber accedido como administrador o responsable.

\textbf{Entradas}:
    \begin{itemize}
    \tightlist
        \item Nombre del usuario.
        \item Apellidos del usuario.
        \item Correo del usuario.
        \item Contraseña.
        \item Teléfono del usuario.
    \end{itemize}

\textbf{Flujo básico}:
    \begin{enumerate}
        \item Se selecciona la opción ``\textbf{\textit{Perfil}}'' del menú lateral.
        
        \item El sistema muestra los datos actuales en los campos correspondientes, excepto la contraseña. Si se trata de un responsable, muestra un listado con los centros y departamentos bajo su responsabilidad.
            
        \item El administrador/responsable puede \textbf{modificar} sus datos.
            \begin{itemize}
            \tightlist
                \item Se guardan los cambios en el botón ``\textbf{\textit{Guardar}}''. 
                    \begin{itemize}
                    \tightlist
                        \item El sistema comprobará que el correo tiene el formato correcto, mostrando un mensaje de error bajo el campo en caso contrario.
                        
                        \item El sistema comprobará que la contraseña tiene como mínimo 6 caracteres incluyendo dígitos y letras minúsculas y mayúsculas, mostrando un mensaje de error ``\textit{Debe contener 6 o más caracteres, incluyendo dígitos, letras minúsculas y mayúsculas}'' bajo el campo en caso contrario. La contraseña se almacena en la base de datos codificada (\textit{hasheada}).
                        
                        \item El sistema comprobará que el teléfono contiene 9 dígitos, mostrando un mensaje de error ``\textit{Debe tener 9 dígitos}'' bajo el campo en caso contrario.
                        
                        \item Se mostrará un mensaje ``\textit{Se han guardado los cambios correctamente}'' si se ha realizado con éxito.
                    \end{itemize}
            \end{itemize}
    \end{enumerate}

\textbf{Salidas}: 
\begin{itemize}
\tightlist
    \item Datos del perfil del usuario, y si se trata de un responsable, listado con los centros y departamentos bajo su responsabilidad.
    
    \item Mensaje de éxito si se ha guardado correctamente.
\end{itemize}

\textbf{Campos de la ventana}:
\tablaSmallSinColores{Campos de la ventana Perfil de Usuario}{llll}{campos_ventana_perfil_usuario}
    {\textbf{Campo} & \textbf{Tipo} & \textbf{Obligatorio} & \textbf{Descripción}\\}{
        \textbf{Nombre} & Texto & Sí & Campo de texto \\ \hline
        \textbf{Apellidos} & Texto & Sí & Campo de texto \\ \hline
        \textbf{Correo} & Texto & Sí & Campo de e-mail \\ \hline
        \textbf{Contraseña} & Texto & Sí & Campo de texto \\ \hline
        \textbf{Teléfono} & Numérico & Sí & Campo numérico \\ 
    }
\apendice{Especificación de diseño}

\section{Introducción}
En este anexo se presentan los aspectos relativos al diseño.

\section{Diseño de datos}
Las entidades que conforman la aplicación son:

\begin{itemize}
    \item \textbf{Usuario}: entidad que almacena los datos de los usuarios que pueden acceder a la aplicación (responsables y administrador). Tiene un id, nombre, apellidos, correo, contraseña, teléfono y un rol. Opcionalmente almacena si el usuario está o no bloqueado, es decir, si se puede eliminar. También mantiene un registro con los centros y departamentos de los que es responsable.
    
    \item \textbf{Propietario aula}: entidad que representa la generalización de los centros y departamentos, cuyo discriminante es el tipo (centro o departamento). Tiene un id y un nombre, además de una referencia al id del usuario que es responsable del mismo. También mantiene un registro de las aulas de las que es responsable.
    
    \item \textbf{Centro}: entidad que representa una especificación de \textit{Propietario aula}, y mantiene un registro de las aulas que están ubicadas en él.
    
    \item \textbf{Departamento}: entidad que representa una especificación de \textit{Propietario aula}.
    
    \item \textbf{Aula}: entidad que almacena los datos de las aulas de los centros y departamentos. Tiene un id, nombre, capacidad y número de ordenadores, además de una referencia al id del centro en el que está ubicada y otra referencia al id de su propietario.
    
    \item \textbf{Reserva}: entidad que almacena los datos de las reservas realizadas sobre las aulas. Tiene un id, fecha que se ha reservado, hora de inicio y de fin de la reserva, día de la semana de la reserva (día de la semana de la fecha), motivo de la reserva (por ejemplo examen, curso, reunión, etc.), persona a cargo de la reserva (no tiene por qué ser el usuario logueado) y usuario responsable de realizar la reserva (el usuario logueado), además de una referencia al id del aula que se ha reservado y otra al id del propietario responsable de la reserva (centro o departamento en nombre del que actúa el usuario). 
    
    \item \textbf{Histórico reservas}: entidad que almacena los datos correspondientes a las operaciones que se han realizado sobre las reservas, las creaciones, modificaciones y eliminaciones. Tiene un id, fecha en que se ha realizado la operación, tipo de operación realizada, los campos identificativos de la reserva sobre opera (motivo, fecha, hora de inicio y de fin, lugar y persona a cargo), usuario responsable de realizar la operación (usuario logueado) y propietario en nombre del que actúa el responsable (propietario del aula de la reserva).
    
    \begin{itemize}
        \item Cuando se trata de una modificación se almacena la operación del histórico con los valores tras la misma.
        
        \item Cuando se trata de un borrado se almacena la operación con los valores de la reserva en el momento del borrado (los de su creación o los de la modificación si se ha realizado).
    \end{itemize}
\end{itemize}

\begin{landscape}
\subsection{Diagrama Entidad/Relación}
\imagenAncho{Diagrama_ER}{Diagrama Entidad/Relación}{1.5}
\end{landscape}

\subsection{Diseño Relacional}
A la hora de realizar el diseño relacional se han tenido en cuenta varias consideraciones:
\begin{itemize}
    \item La tabla \textbf{\textit{historico\_reservas}} se trata de una tabla independiente, es decir, no mantiene ninguna relación con las otras tablas, para poder almacenar la información de las operaciones realizadas sobre las reservas a pesar de que éstas, los usuarios, las aulas o los propietarios se borren de la base de datos. 
    
        \begin{itemize}
            \item Esta decisión fuerza a tener que almacenar información aparentemente redundante sobre los datos de las reservas, ya que al no tener establecida la \textit{foreign key} con \textit{reserva} no se tiene acceso a sus datos.
            
            \item Si se hubieran incluido las relaciones con las demás tablas, como por ejemplo con la tabla ``\textit{reserva}'', cuando se produce un borrado de la reserva (algo que puede ser bastante habitual), el ID de la reserva incluido en el registro del histórico quedaría huérfano, produciendo una excepción de violación de dicha clave ajena.
        \end{itemize}
        
    \item La tabla \textbf{\textit{reserva}} se ha desnormalizado ~\cite{8104}. 
        \begin{itemize}
            \item No alcanza la 3FN ya que se ha añadido a la misma un atributo que especifica el usuario responsable de realizar la reserva (el usuario que se ha logeado).
            \item Al añadir este atributo, aparece la dependencia funcional \textit{aula --> usuario\_responsable}, ya que la reserva determina el aula, el aula determina el propietario, y el propietario determina el responsable. En la dependencia, el determinante \textit{aula} no es una clave de la tabla \textit{reserva} (i.e., no está en FNBC), ni si quiera es parte de su clave primaria (i.e., tampoco está en 3FN), pero sí está en 2FN ya que la clave de esta tabla \textit{reserva} es una clave simple.
            \item Al igual que en el caso anterior, se ha tomado esta decisión para evitar problemas cuando se elimine un usuario.
        \end{itemize}
\end{itemize}

\begin{landscape}
\imagenAncho{Diagrama_Relacional}{Diagrama Relacional}{1.5}
\end{landscape}

\section{Diseño arquitectónico}
\subsection{Diseño de paquetes}
En lo relacionado al diseño de paquetes del proyecto se pueden destacar que varios de ellos que interactúan con componentes externos.

\begin{itemize}
    \item \textit{\textbf{app.security}}: es el paquete que contiene lo relativo a las configuraciones de seguridad. 
    
    \begin{itemize}
        \item Interactúa con Vaadin, implementando \textit{VaadinServiceInitListener}, para conectar Spring Security a su sistema de navegación.
        
        \item Define un interfaz funcional \textit{CurrentUser} relacionándose con la entidad \textit{Usuario} definida en el \textit{backend} para obtener el usuario logeado.
        
        \item Extiende de \textit{HttpSessionRequestCache} para poder realizar un seguimiento de la URL a la que se intenta acceder.
        
        \item Extiende de \textit{WebSecurityConfigurerAdapter} para crear la clase de configuración de seguridad, incluyendo anotaciones de Spring como \textit{@EnableWebSecurity} (para aplicar indicar que se aplica Spring Security) y \textit{@Configuration} (para indicar que se debe usar dicha clase para configurar la seguridad).
        
        \item Implementa \textit{UserDetailsService} del framework de Spring, para buscar los usuarios (de la entidad \textit{Usuario} del backend) a partir del correo proporcionado en el login. 
    \end{itemize}
    
    \item \textit{\textbf{backend}}: paquete que define las entidades, servicios y repositorios para conectar la aplicación con la base de datos a través de Spring JPA~\cite{pagina_spring_jpa}.
    
    \item \textit{\textbf{exceptions}}: paquete que define las excepciones propias de la aplicaicón, extendiendo de, por ejemplo, \textit{DataIntegrityViolationException} o \textit{RuntimeException}.
    
    \item \textit{\textbf{ui}}: paquete que define las ventanas de la aplicación, interactuando con el \textit{backend}, el \textit{frontend}, el paquete de seguridad \textit{app.security} y de excepciones \textit{exceptions}, y los componentes específicos de Vaadin.
\end{itemize}


\begin{landscape}
\imagenAncho{Diagrama_Paquetes}{Diagrama de Paquetes}{1.5}
\end{landscape}
\apendice{Documentación técnica de programación}

\section{Introducción}
En este anexo se explica la estructura de la que dispone el proyecto, se incluye un manual del programador para poder instalar el entorno y continuar con el desarrollo y por último la compilación, instalación y ejecución del proyecto. Además, también se incluyen las pruebas realizadas.

\section{Estructura de directorios}
Las aplicaciones web desarrolladas con Vaadin son \textit{full stack}, incluyendo el código del lado del cliente (\textit{frontend}) y del lado del servidor (\textit{backend}) en el mismo proyecto.

La estructura que sigue el repositorio es la siguiente:
\begin{itemize}
	\item \textbf{/}: archivos de configuración de Vaadin y Spring Boot, fichero README y LICENSE.
	
	\item \textbf{/.vaadin/designer}: archivo de configuración para el diseño de ventanas con JavaScript.
	
	\item \textbf{/Documentación}: documentación del proyecto.
	\begin{itemize}
	\tightlist
    	\item \textbf{/Diagramas}: archivos con los diagramas creados en StarUML.
    	
    	\item \textbf{/Versiones}: histórico con los documentos de las especificaciones de requisitos y prototipos de ventanas de las diferentes versiones anteriores a la actual.
    	
    	\begin{itemize}
    	\tightlist
        	\item \textbf{/Versión 00}: documentos con las especificaciones de requisitos y prototipos de ventanas de la Versión 0 (versión de partida antes de comenzar con el desarrollo).
        	
        	\item \textbf{/Versión 1}: documentos con las especificaciones de requisitos y prototipos de ventanas de la Versión 1, tras la última revisión del tutor. 
        \end{itemize}
    \end{itemize}
    
    \item \textbf{/frontend}: código del lado de cliente.
        \begin{itemize}
    	\tightlist
        	\item \textbf{/src/views/errors}: archivo JavaScript que define el diseño de la ventana que se muestra ante un acceso denegado.
        	
        	\item \textbf{/styles}: hojas de estilo (archivos CSS) para el diseño de las ventanas.
        \end{itemize}
        
    \item \textbf{/sql}: script SQL para la generación de la base de datos de PostgreSQL.
    
    \item \textbf{/src}: código del lado del servidor (\textit{backend}).
        \begin{itemize}
    	\tightlist
        	\item \textbf{/main}: código de generación de las ventanas.
        	
        	    \begin{itemize}
                	\item \textbf{/java/gestionaulasinformatica}: 
                	
                    	\begin{itemize}
                        	\item \textbf{/app/security}: código de configuración de la seguridad de la aplicación con Spring Security.
                        	
                        	\item \textbf{/backend}: código de datos, entidades, servicios, repositorios y especificaciones.
                	
                        	\item \textbf{/backend/data}: código que define tipos de datos, como los roles de usuario, el tipo de operaciones sobre las reservas o el tipo de propietario de aulas.
                        	
                        	\item \textbf{/backend/entity}: código de las entidades JPA correspondientes a las tablas de la base de datos.
                        	
                        	\item \textbf{/backend/respository}: repositorios JPA para cada entidad.
                        	
                        	\item \textbf{/backend/service}: servicios JPA que implementan los repositorios para cada entidad.
                        	
                        	\item \textbf{/backend/specification}: especificaciones de ciertas entidades para realizar filtrado de datos.
                        	
                        	\item \textbf{/exceptions}: código de excepciones propias.
                        	
                        	\item \textbf{/ui}: código de las ventanas, archivo que define la plantilla de las ventanas, archivo con funciones comunes y enumeración con alguno de los mensajes que se muestran.
                            
                            \item \textbf{/ui/views}: código de las ventanas.
                        	
                        \end{itemize}
                	
                	\item \textbf{/resources}: archivo con las propiedades para la configuración de la aplicación.
                \end{itemize}
                
        	\item \textbf{/test}: código de generación de los tests.
        	    \begin{itemize}
            	\tightlist
                	\item \textbf{/java/gestionaulasinformatica/backend}: código de los tests. La carpeta debe tener la misma ruta que aquella en la que se encuentra el código que se está testeando.
                	
                	\item \textbf{/resources}: archivo con las propiedades para la configuración de la aplicación para los test.
                \end{itemize}
        \end{itemize}
	
\end{itemize}

\section{Manual del programador}
En este manual se detalla la instalación de los programas necesarios para continuar con el desarrollo de la aplicación.

\subsection{Instalación de Java}
Es necesario trabajar con la versión 8 de Java, por lo tanto hay que tener instalado el JDK 8.

En el desarrollo del proyecto se ha trabajado con la \textbf{versión 1.8.0\_251}, aunque es recomendable descargar la última que esté disponible.

Para instalar JDK 8 hay que acceder a la \href{https://www.oracle.com/java/technologies/javase-downloads.html}{página de descargas de Java SE}, buscar ``\textbf{Java SE 8u...}'' y descargar el JDK.

\imagenAncho{Descarga_JDK_8}{Descarga de JDK 8}{0.9}

Al ejecutar el instalador debemos leer y aceptar la licencia de uso de Oracle, seleccionar el sistema operativo y arquitectura del ordenador en que se va a instalar y seguir los pasos del asistente.

\subsection{Instalación de Eclipse}
El siguiente paso es instalar un IDE de Java, en este caso se ha utilizado \textbf{Eclipse IDE for Eclipse Committers}.

Para ello hay que acceder a la \href{https://www.eclipse.org/downloads/}{página de descargas de Eclipse} y descargar la opción ``\textit{Get Eclipse IDE\dots}''.

\imagenAncho{Descarga_Eclipse_IDE}{Descarga de Eclipse IDE}{0.7}

Al ejecutar el instalador hay que buscar y seleccionar ``\textbf{\textit{Eclipse IDE for Eclipse Commiters}}'' y después en el apartado ``\textbf{\textit{Java 1.8+ VM}}'' seleccionar el JDK 8 que se ha instalado previamente.

\imagenAncho{Instalador_Eclipse}{Instalador de Eclipse}{0.55}
\imagenAncho{Instalador_Eclipse_JDK}{Instalador de Eclipse - Selección de Java}{0.6}

\subsection{Instalación del \textit{plugin} de Vaadin para Eclipse}
Tras instalar Eclipse, lo siguiente que hay que hacer es añadir el \textit{plugin} de Vaadin para Eclipse.~\footnote{El \textit{plugin} es indepenpendiente a la versión de Vaadin utilizada para el desarrollo, esta se indica directamente en el archivo \textit{pom.xml}.}

Para ello hay que abrir Eclipse y en la barra de herramientas seleccionar ``\textbf{\textit{Help/Eclipse Marketplace\dots}}'', en la nueva ventana hay que buscar ``\textbf{\textit{vaadin}}'', pulsar ``\textbf{\textit{Go}}'' y una vez aparece ``\textbf{\textit{Vaadin Plugin for Eclipse}}'' pulsar en ``\textbf{\textit{Install}}''. ~\footnote{Una vez se ha instalado, como en mi caso, aparece como ``\textit{Installed}''.}

\imagenAncho{Eclipse_marketplace}{Eclipse Marketplace}{0.5}
\imagenAncho{Eclipse_Marketplace_Plugin_Vaadin}{Eclipse Marketplace - Plugin de Vaadin}{0.8}

\subsection{Instalación de PostgreSQL y pgAdmin}
Para poder administrar la base de datos hay que instalar PostgreSQL junto a su herramienta de gestión pgAdmin.

Desde la \href{https://www.enterprisedb.com/downloads/postgres-postgresql-downloads}{página de descarga de PostgreSQL}, se puede descargar clicando en``\textbf{\textit{Download}}'' de la fila correspondiente a la \textbf{versión 12.3 para Windows}. Lo siguiente es ejecutar el instalador siguiendo los pasos del asistente. 

\imagen{Descarga_PostrgeSQL_pgAdmin}{Descarga de PostgreSQL y pgAdmin}

\subsection{Instalación de Heroku CLI}
Por último, para poder desplegar la aplicación y manejar la base de datos del despliegue más fácilmente, hay que instalar el cliente de Heroku, el cual permite trabajar desde la consola de comandos.

Para ello hay que acceder a la \href{https://devcenter.heroku.com/articles/heroku-cli#download-and-install}{página de descarga de Heroku CLI}, elegir el instalador de la arquitectura Windows del ordenador en el que se esté trabajando y ejecutar el instalador siguiendo el asistente.

\imagen{Descarga_Heroku_CLI}{Descarga de Heroku CLI}

\section{Compilación, instalación y ejecución del proyecto}
En el desarrollo del proyecto, como se ha comentado anteriormente, se ha utilizado Eclipse, por lo que en esta guía se explica como llevar a cabo el proceso en este IDE.

\subsection{Descarga del código fuente}
El código fuente del proyecto se encuentra hospedado en GitHub, por lo que se puede descargar fácilmente desde el mismo repositorio.

Para ello hay que acceder al \href{https://github.com/lcs1001/gestion-aulas-informatica}{repositorio del proyecto} y clicar en el botón ``\textbf{\textit{Code}}'' que aparece en verde. 

El proyecto se puede descargar con Git, copiando la URL, abrirlo con GitHub Desktop, abrirlo con Visual Studio o descargar directamente el archivo ZIP. En este caso se va a explicar la instalación mediante el ZIP, por lo que hay que clicar en ``\textbf{\textit{Download ZIP}}''.

\imagen{Descarga_codigo_fuente}{Descarga del código fuente}

Lo siguiente es importar el proyecto a Eclipse, para ello hay que descomprimir el ZIP descargado, abrir el IDE y en la barra de herramientas seleccionar \textbf{\textit{File/Import\dots}}.

\imagenAncho{Eclipse_importar}{Importar el código fuente a Eclipse (1)}{0.5}

En la ventana que se abre hay que seleccionar la opción ``\textbf{\textit{Projects from Folder or Archieve}}'' y clicar en ``\textbf{\textit{Next >}}''.

\imagen{Eclipse_importar_archivo}{Importar el código fuente a Eclipse (2)}

En la siguiente ventana hay que clicar en ``\textbf{\textit{Directory\dots}}'', seleccionar la carpeta descomprimida con nombre ``\textbf{gestion-aulas-informatica-master}''\ y clicar en ``\textbf{\textit{Finish}}'' una vez se ha cargado.

\imagen{Eclipse_importar_archivo_2}{Importar el código fuente a Eclipse (3)}


\subsection{Configuración de la base de datos}
El siguiente paso es generar y configurar la base de datos a partir del script SQL para poder administrarla desde pgAdmin.

Lo primero que hay que hacer es crear la base de datos, para ello se abre pgAdmin y tras introducir la contraseña aparecerá en la parte derecha el desplegable ``\textit{Servers}'', hay que abrir los desplegables hasta llegar al apartado ``\textit{Databases}'' y hacer clic derecho ``\textbf{\textit{Create/Database\dots}}''. 

\imagenAncho{pgAdmin_crear_bd}{Creación de la base de datos ``gestionaulasinformatica'' (1)}{0.5}

En la ventana que se abre hay que introducir en el apartado ``\textit{Database}'' el nombre ``gestionaulasinformatica'' y hacer clic en ``\textbf{\textit{Save}}''.

\imagenAncho{pgAdmin_crear_bd_2}{Creación de la base de datos ``gestionaulasinformatica'' (2)}{0.5}

Tras crear la base de datos hay que acceder a la carpeta donde se ha instalado PostgreSQL y entrar en la carpeta ``\textbf\textit{{bin}}'', una vez ahí se escribe en la barra en que se muestra la ruta la palabra ``\textbf{\textit{cmd}}'' para abrir la consola de comandos en esta ruta.

\imagen{PostgreSQL_generar_bd}{Acceder a la consola de comandos para generar la base de datos}

Una vez abierta la consola de comandos, se escribe el comando ``\textbf{\textit{psql --username postgres --dbname gestionaulasinformatica -f}}'' seguido de la ruta en que se haya descargado el proyecto y la ruta al script. por ejemplo si se ha descargado en la carpeta ``\textit{Descargas}'' el comando completo sería ``\textit{psql --username postgres --dbname gestionaulasinformatica -f C:\textbackslash{}Users\textbackslash{}<usuario>\textbackslash{}Downloads\textbackslash{}gestion-aulas-informatica-master\textbackslash{}sql\textbackslash{} Script\_GestionAulasInformatica}''.

Tras configurar la base de datos, hay que, como mínimo, insertar un registro con los datos del administrador en la tabla ``\textit{usuario}'' para poder acceder a la aplicación y empezar a crear el resto de usuarios (responsables) y los demás datos de centros, departamentos y aulas.

Para ello hay que acceder a la tabla desde ``\textbf{\textit{gestionaulasinformatica/ Schemas/Tables}}'' y en ``\textbf{\textit{usuario}}'' hacer click derecho ``\textbf{\textit{Scripts/INSERT Script}}'', por lo que se abrirá una ventana con el script de inserción.

La primera vez que se inserta un registro de usuario hay que introducir en la línea anterior al \textit{insert} el comando ``\textbf{\textit{CREATE EXTENSION pgcrypto;}}'', para poder insertar las contraseñas encriptadas en la base de datos.

El siguiente paso es insertar los valores del usuario administrador:
\begin{itemize}
    \item En el campo \textit{id\_usuario} se pone el valor ``\textit{default}''. ~\footnote{Los campos autonuméricos deben insertarse como ``\textit{default}'' para seguir la secuencia generada.}
    
    \item En el campo \textit{contraseña} se utiliza el comando ``\textit{crypt('\textbf{contraseña}', gen\_salt('bf'))}'' para almacenarla encriptada.
    
    \item En el campo \textit{rol} se pone el valor ``\textit{ADMIN'}'.
    
    \item En el campo \textit{bloqueado} se pone el valor ``\textit{true}''. ~\footnote{Con el valor \textit{true} en el campo ``\textit{bloqueado}'' se indica que el usuario no puede ser eliminado.}
\end{itemize}

Un ejemplo de usuario administrador sería el siguiente:

\imagenAncho{pgAdmin_ejemplo_administrador}{Ejemplo de inserción de un registro en la tabla ``\textit{usuario}'' del usuario administrador}{1}

Para insertar el registro en la tabla se puede pulsar ``\textit{F5}'' o clicar en el botón con el icono \textit{play} de la barra de herramientas (véase el cuadro rojo de la imagen anterior).

El último paso es especificar el usuario, en este caso ``\textit{postgre}'', y la contraseña de acceso a la base de datos en el archivo \textbf{\textit{application\_properties}} (que se encuentra en la carpeta \textit{gestion-aulas-informatica-master/src/main/resources} del proyecto) desde Eclipse. Esto es necesario para poder ejecutar la aplicación en local.

\imagen{Eclipse_configurar_bd_local}{Configuración del archivo \textit{application\_properties} para especificar usuario y contraseña de la base de datos}

\subsection{Ejecución del proyecto}
Para ejecutar el proyecto hay que acceder a la carpeta ``\textbf{\textit{gestion-aulas-informatica-master/src/main/java/gestionaulasinformatica}}'' y en el archivo ``\textbf{\textit{Application.java}}'' hacer click derecho y seleccionar ``\textbf{\textit{Run As/Java Application}}''.

\imagen{Eclipse_ejecutar}{Ejecución del proyecto}

\subsection{Despliegue en Heroku}
Lo primero que se debe hacer, si no se dispone de una cuenta, es darse de alta desde la \href{https://www.heroku.com/}{página de Heroku}.

\subsubsection{Creación de la app}
El siguiente paso es crear una nueva \textit{app} clicando en ``\textbf{\textit{Create new app}}'' en la ventana principal.

\imagenAncho{Heroku_creacion_app}{Creación de la \textit{app} en Heroku(1)}{0.9}

En la siguiente ventana hay que indicar el nombre, en este caso ``\textit{gestionaulasinformatica}'', y la zona en la que se va a desplegar \textit{Europa} y clicar en ``\textbf{\textit{Create app}}''.

\imagenAncho{Heroku_creacion_app_2}{Creación de una nueva app en Heroku (2)}{0.9}

No es necesario realizar ninguna configuración más, ya que por defecto está marcada la opción de despliegue con el cliente heroku (\textit{Heroku Git}).

\imagenAncho{Heroku_app_forma_despliegue}{Configuración de despliegue de la \textit{app} en Heroku}{1}

\subsubsection{Configuración de la base de datos PostgreSQL}
Después crear la \textit{app}, hay que añadir la extensión de PostgreSQL para manejar la base de datos del despliegue. 

Para ello se puede seguir el \href{https://devcenter.heroku.com/articles/heroku-postgresql#provisioning-heroku-postgres}{tutorial de la página oficial de Heroku}, añadiendo al final del comando ``\textit{-a gestionaulasinformatica}'' para indicar que se quiere incluir en la \textit{app} creada anteriormente. ~\footnote{En el tutorial se muestra directamente la configuración con el plan gratuito, que es el que se ha utilizado.}

Tras la configuración, aparecerá en la ventana principal la extensión añadida:

\imagen{Heroku_complemento_PostrgreSQL}{Extensión de PostgreSQL para Heroku (1)}

Por último, queda conectarse al servidor de la base de datos que ha creado Heroku a pgAdmin para poder administrarla. ~\footnote{Para ello se ha seguido el tutorial que se encuentra en ~\cite{importar_bd_postgresql_heroku}.}

Hay que acceder a la extensión de PostgreSQL al apartado ``\textbf{\textit{Settings}}'' y clicar en ``\textbf{\textit{View credentials\dots}}'' para ver las credenciales que se necesitan para la conexión.

\imagenAncho{Heroku_complemento_PostrgreSQL_credenciales_1}{Acceso a las credenciales de la extensión PostgreSQL de Heroku (1)}{1}

\imagenAncho{Heroku_complemento_PostrgreSQL_credenciales_2}{Acceso a las credenciales de la extensión PostgreSQL de Heroku (2)}{0.5}

Una vez se tienen las credenciales, hay que acceder a pgAdmin y crear un nuevo servidor haciendo click derecho en ``\textbf{\textit{Servers}}'' y seleccionando ``\textbf{\textit{Create/Server\dots{}}}''.

\begin{itemize}
    \item En la pestaña ``\textbf{\textit{General}}'' se introduce el nombre que se le va a dar al servidor en ``\textbf{\textit{Name}}''.
    
    \item En la pestaña ``\textbf{\textit{Connection}}'' se introduce:
        \begin{itemize}
            \item En ``\textbf{\textit{Host}}'' el campo ``\textit{host}'' de las credenciales.
            
            \item En ``\textbf{\textit{Port}}'' se deja el ``5432''.
            
            \item En ``\textbf{\textit{Maintenance database}}'' el campo ``\textit{database}'' de las credenciales.
            
            \item En ``\textbf{\textit{Username}}'' el campo ``\textit{user}'' de las credenciales.
            
            \item En ``\textbf{\textit{Password}}'' el campo ``\textit{password}'' de las credenciales.
        \end{itemize}
        
    \item En la pestaña ``\textbf{\textit{SSL}}'' se selecciona el ``\textbf{\textit{SSL mode}}'' como ``\textbf{\textit{Require}}''.
    
    \item En la pestaña ``\textbf{\textit{Advanced}}'', en ``\textbf{\textit{DB restriction}}'' se introduce el campo ``\textit{Maintenance database}'' de las credenciales. ~\footnote{De esta forma se evita que aparezcan todas las bases de datos del servidor de PostgreSQL de Heroku.}
\end{itemize}

Al finalizar hay que clicar en ``\textbf{\textit{Save}}'' y ya aparecerá la base de datos en la lista.

El último paso tras importar la base de datos es generar las tablas con el script SQL al igual que se ha hecho en local. Para ello se abre la consola de comandos y se escribe el comando ``\textbf{\textit{heroku pg:psql --app gestionaulasinformatica <}}'' seguido de la ruta en que se haya descargado el proyecto y la ruta al script. por ejemplo si se ha descargado en la carpeta ``\textit{Descargas}'' el comando completo sería ``\textit{heroku pg:psql --app gestionaulasinformatica < C:\textbackslash{}Users\textbackslash{}<usuario>\textbackslash{}Downloads\textbackslash{}gestion-aulas-informatica-master\textbackslash{}sql\textbackslash{} Script\_GestionAulasInformatica}''.

\subsection{Despliegue}
Después de haber creado la aplicación en Heroku y configurado la base de datos, queda empaquetar el proyecto y desplegar la \textit{app}, para ello hay que acceder a la carpeta en la que se encuentra el proyecto (en el administrador de archivos) y escribir la palabra ``\textbf{\textit{cmd}}'' en la barra en que se muestra la ruta la  para abrir la consola de comandos.

Para empaquetar el proyecto como \textit{.JAR} hay que ejecutar el comando ``\textit{mvn package -Pproduction}'' (se indica que se está empaquetando en modo producción). Si todo ha ido bien, aparecerá lo siguiente en la consola:

\imagenAncho{mvn_package}{Despliegue en Heroku - Empaquetado de la aplicación en .JAR desde \textit{cmd}}{1}

El siguiente y último paso es el despliegue, para ello se ejecuta el comando ``\textit{heroku deploy:jar target/gestion-aulas-informatica-2.0-SNAPSHOT.jar --app gestionaulasinformatica}''. Si se ha realizado correctamente, aparecerá lo siguiente en la consola:

\imagenAncho{heroku_deploy}{Despliegue en Heroku - Despliegue desde \textit{cmd}}{1}

Para abrir la aplicación una vez se ha desplegado se puede ejecutar el comando ``\textit{heroku open --app gestionaulasinformatica}'' en la consola o acceder directamente a través de la URL \url{https://gestionaulasinformatica.herokuapp.com/}.

\subsection{Continuación del desarrollo tras el despliegue}
Al empaquetar el proyecto como producción, es necesario recompilar el proyecto ya que muchos archivos desaparecen al activar este modo, para ello se ejecuta el comando ``\textit{mvn compile}'' en la consola de comandos (en la carpeta en la que se encuentra el proyecto). Si se ha compilado correctamente, aparecerá un mensaje similar al del empaquetamiento:

\imagenAncho{mvn_compile}{Continuación del desarrollo tras el despliegue - Recompilación del proyecto}{1}

\section{Pruebas del sistema}
Para poder verificar el funcionamiento de ciertas partes de la aplicación, y con el fin de no tener que realizar las pruebas manualmente cada vez que cambiaba el código, se han desarrollado una serie de tests de integración con la ayuda de Spring Boot y su dependencia para la parte de \textit{testing}~\cite{pagina_spring_boot_testing}, añadiendo la anotación ``\textit{@DataJpaTest}'' enla declaración de la clase. En el fondo estos test trabajan internamente con JUnit 5~\cite{pagina_junit}.

En cada test se generan los datos y se revierten una vez ejecutados, y además se imprimen mensajes en la consola de ejecución con el test que se está realizando y el resultado obtenido.

Para ejecutar todos los tests hay que buscar la carpeta ``\textit{src/test/java}'' en el proyecto de Eclipse y hacer click derecho ``\textbf{\textit{Run As/JUnit Test}}''. O si se quiere ejecutar uno a uno por separado, en la misma carpeta se hace click derecho ``\textbf{\textit{Run As/JUnit Test}}'' sobre la clase que se quiera probar.

\imagenAncho{Test_ejecutar}{Pruebas del sistema - Ejecución de los test en Eclipse}{1}

Al ejecutar los tests, Eclipse cambiará a la vista de JUnit, mostrando en una pestaña el resultado de los mismos. Si se han ejecutado correctamente aparecerán todos con un check en verde:

\imagenAncho{Test_ejecutar_correctamente}{Pruebas del sistema - Test ejecutados con éxito}{0.5}

\subsection{Pruebas para la Consulta de Disponibilidad de Aulas}
Para \textit{testear} la consulta de disponibilidad de aulas se instancian 2 usuarios con rol de responsables, 1 centro y 1 departamento, 2 aulas (ambas bajo responsabilidad del departamento) y 3 reservas (simulando una reserva por rango de fechas).

Se han creado 3 tests:
\begin{itemize}
    \item En el primero se comprueba la disponibilidad en una fecha y horas concretas.
    \item En el segundo se comprueba lo mismo, añadiendo el filtro ``\textit{Día de la semana}''.
    \item En el último se comprueba la disponibilidad del aula en un rango de fechas y unas horas concretas.
\end{itemize}

\imagenAncho{Test_consulta_aulas}{Pruebas del sistema - Test para la Consulta de Disponibilidad de Aulas}{0.5}

La salida en consola de uno de ellos se vería así:

\imagenAncho{Test_consulta_aulas_consola}{Pruebas del sistema - Test para la Consulta de Disponibilidad de Aulas - Salida en consola}{0.9}

\subsection{Pruebas para la entidad PropietarioAula}
Para \textit{testear} la entidad PropietarioAula se han instanciado un usuario con rol de responsable y un centro.

Se han creado 2 tests:
\begin{itemize}
    \item En el primero se comprueba el funcionamiento del método \textit{findById}.
    \item En el segundo se comprueba el funcionamiento del método \textit{delete}.
\end{itemize}

\imagenAncho{Test_propietario}{Pruebas del sistema - Test para la entidad PropietarioAula}{0.4}

La salida en consola de uno de ellos se vería así:

\imagenAncho{Test_propietario_consola}{Pruebas del sistema - Test para la entidad PropietarioAula - Salida en consola}{1}

\subsection{Pruebas para la entidad Usuario}
Para \textit{testear} la entidad Usuario se ha instanciado un usuario con rol de responsable.

Se ha creado un único test para comprobar el método \textit{findById}.

\imagenAncho{Test_usuario}{Pruebas del sistema - Test para la entidad Usuario}{0.4}

La salida en consola se vería así:

\imagenAncho{Test_usuario_consola}{Pruebas del sistema - Test para la entidad Usuario - Salida en consola}{0.6}

\subsection{Pruebas manuales}
Se ha reforzado la comprobación del correcto funcionamiento de toda la aplicación con pruebas manuales.

Pruebas realizadas para el \textbf{Login}:
\begin{itemize}
    \item Introducir el correo y/o la contraseña incorrectos y comprobar que limpia los campos.
    \item Introducir un correo y contraseñas cualquiera (no registradas) y comprobar que limpia los campos.
    \item Introducir los datos del administrador y comprobar que únicamente se muestran en el menú lateral las ventanas accesibles por el mismo. 
    \item Introducir los datos de un usuario con rol responsable y comprobar que únicamente se muestran en el menú lateral las ventanas accesibles por el mismo.
    \item Habiendo accedido como administrador, comprobar que si se escribe en la URL la ruta de una de las ventanas del usuario con rol de responsable, aparece la pantalla de acceso denegado. Y viceversa.
    \item Acceder a la consulta de reservas y disponibilidad de aulas a través del enlace de la ventana de login, sin haberse logueado, y comprobar que sólo se muestran dichas ventanas en el menú lateral. 
    \item Comprobar que si se intenta acceder a una de las ventanas del administrador o del usuario con rol de responsable, sin haberse logueado, se redirige a la ventana de login.
\end{itemize}

Pruebas realizadas para la \textbf{Consulta de Reservas}:
\begin{itemize}
    \item Dejar el campo \textit{Centro/Departamento} vacío y comprobar que se muestra la notificación de error.
    \item Comprobar que en el listado se muestran solo las reservas que cumplen con los filtros aplicados, mostrándose una notificación de que no hay ninguna si no coinciden.
    \item Comprobar que al clicar en \textit{Limpiar filtros} se vacían los campos de entrada y se pone en \textit{Fecha desde} la fecha actual.
\end{itemize}

Pruebas realizadas para la \textbf{Consulta de Disponibilidad de aulas}:
\begin{itemize}
    \item Dejar el campo \textit{Centro/Departamento} vacío y comprobar que se muestra la notificación de error.
    \item Rellenar el campo \textit{Fecha desde} y comprobar que se muestra la notificación de error.
    \item Comprobar que en el listado se muestran solo las aulas disponibles que cumplen con los filtros aplicados, mostrándose una notificación de que no hay ninguna si no coinciden.
    \item Comprobar que al clicar en \textit{Limpiar filtros} se vacían los campos de entrada.
\end{itemize}

Pruebas realizadas para el \textbf{Mantenimiento de Centros y Departamentos}:
\begin{itemize}
    \item Crear un centro y comprobar que aparece en el listado.
    \item Modificar los datos del centro y comprobar que aparecen actualizados en el listado.
    \item Eliminarlo sin haberle asignado aulas.
    \item Crear un centro, asignarle un aula y eliminarlo verificando que en el mensaje de confirmación se indica que tiene aulas bajo su responsabilidad y que se van a eliminar si se confirma. Comprobando también que se eliminan las aulas que están asociadas a él.
    \item Comprobar que el filtrado por nombre funciona.
\end{itemize}

Pruebas realizadas para el \textbf{Mantenimiento de Aulas}:
\begin{itemize}
    \item Crear un aula y comprobar que aparece en el listado.
    \item Modificar los datos del aula y comprobar que aparecen actualizados.
    \item Eliminarlo sin haberle asignado reservas.
    \item Crear un aula, asignarle una reserva y eliminarla verificando que en el mensaje de confirmación se indica que tiene reservas asignadas y que se van a eliminar si se confirma. Comprobando también que se eliminan las reservas que están asociadas al aula.
\end{itemize}

Pruebas realizadas para el \textbf{Mantenimiento de Usuarios}:
\begin{itemize}
    \item Crear un usuario y comprobar que aparece en el listado.
    \item Modificar los datos del usuario y comprobar que aparecen actualizados.
    \item Eliminarlo sin haberle asignado centros o departamentos.
    \item Crear un usuario, asignarle un centro y eliminarla verificando que informa de que primero debe reasignar los centros o departamentos bajo su responsabilidad. Reasignar el centro a otro responsable y comprobar que permite borrarlo.
\end{itemize}

Pruebas realizadas para el \textbf{Histórico de Reservas}:
\begin{itemize}
    \item Filtrar por \textit{Fecha desde} y comprobar que en el listado sólo aparecen las operaciones realizadas a partir de dicha fecha.
    \item Filtrar por \textit{Fecha hasta} y comprobar que en el listado sólo aparecen las operaciones realizadas hasta dicha fecha.
    \item Filtrar por \textit{Fecha desde} - \textit{Fecha hasta} y comprobar que en el listado sólo aparecen las operaciones realizadas entre ese rango de fechas.
    \item Comprobar que al clicar en \textit{Limpiar filtros} se vacían los campos de entrada.
\end{itemize}

Pruebas realizadas para el \textbf{Perfil}:
\begin{itemize}
    \item Cambiar algún campo y comprobar que al guardar se actualizan los datos correctamente.
    \item Introducir un correo con formato incorrecto y comprobar que aparece el mensaje de error bajo el campo.
    \item Introducir una contraseña con formato incorrecto y comprobar que aparece el mensaje de error bajo el campo.
    \item Introducir un teléfono con formato incorrecto y comprobar que aparece el mensaje de error bajo el campo.
\end{itemize}

Pruebas realizadas para la \textbf{Reserva de Aulas}:
\begin{itemize}
    \item Hacer una reserva y comprobar que se ha registrado correctamente accediendo a la Gestión de Reservas y viendo si aparece en el listado.
    \item Hacer una reserva del mismo aula que colisione con la anterior, es decir, que coincida en fecha y parte del horario, y comprobar que aparece la notificación de error de que el aula no está disponible.
    \item Hacer lo mismo con las reservas de rangos de fechas, incluyendo el día de la semana y sin incluirlo.
\end{itemize}

Pruebas realizadas para la \textbf{Gestión de Reservas}:
\begin{itemize}
    \item Filtrar las reservas que se muestran y comprobar que funciona correctamente, del estilo a las realizadas en la Consulta de Reservas.
    \item Comprobar que al clicar en \textit{Limpiar filtros} se vacían los campos de entrada y se pone en \textit{Fecha desde} la fecha actual.
    \item Sin seleccionar ninguna reserva clicar en \textit{Modificar} y comprobar que se muestra la notificación de error de que no se ha incluido ninguna, y lo mismo clicando en \textit{Eliminar}.
    \item Seleccionar varias reservas y clicando en \textit{Modificar}, comprobar que se muestra la notificación de error de que no se puede modificar más de una reserva a la vez.
    \item Seleccionar una reserva, clicar en \textit{Modificar} y comprobar que se actualizan los campos correctamente en el listado. Comprobando también que se muestra la notificación de error si el aula no está disponible.
    \item Seleccionar una o varias reservas y comprobar que se eliminan correctamente, mostrando primero un mensaje de confirmación.
\end{itemize}
\apendice{Documentación de usuario}

\section{Introducción}
En este anexo se indican los requisitos mínimos para que el usuario pueda hacer uso de la aplicación y el manual de uso de la misma.

\section{Requisitos de usuarios}
Como se trata de una aplicación web, el único requisito con el que debe cumplir el usuario es disponer de un dispositivo con acceso a Internet, ya que la aplicación se puede usar tanto en teléfonos móviles, \textit{tablets} y ordenadores. Aunque por comodidad, es recomendable acceder a la misma desde un ordenador.

\section{Instalación}
No es necesario llevar a cabo ninguna instalación, ya que se trata de una aplicación online, a la que se puede acceder a través del enlace \url{https://gestionaulasinformatica.herokuapp.com}.

\section{Manual del usuario}
En este apartado se describe el uso básico de la aplicación.

\subsection{Acceso como usuario visitante}
Cualquier usuario puede acceder a la consulta de reservas y a la consulta de disponibilidad de aulas sin necesidad de estar registrado en el sistema (sin \textit{login}).

Para ello hay que clicar en la ventana de inicio en el enlace ``\textbf{\textit{Acceso consulta de reservas y disponibilidad de aulas}}'' que se encuentra bajo el formulario de \textit{login}.

\imagenAncho{Login_acceso_consulta}{Acceso a la consulta de reservas y disponibilidad de aulas como visitante}{0.5}

El enlace redirige a la ventana de consulta de reservas, clicando en el botón del menú lateral (que se encuentra en la parte superior izquierda) se puede acceder a la consulta de disponibilidad de aulas y clicando en ``\textit{Iniciar sesión}'' se vuelve a la ventana de login.

\imagenAncho{Visitante_consulta_reservas}{Ventana Consulta de Reservas con acceso como visitante}{0.9}

\subsection{Login}
Para acceder al sistema hay que estar dado de alta previamente como usuario responsable o administrador.
Se introduce el correo de usuario en el campo ``\textbf{\textit{Username}}'' y la contraseña en el campo ``\textbf{\textit{Password}}'' y se hace click en el botón ``\textbf{\textit{Log in}}''.

\begin{itemize}
    \item Si la contraseña introducida es incorrecta se limpiarán los campos introducidos.
    \item Si el usuario no recuerda la contraseña puede hacer click en ``\textbf{\textit{Forgot password}}'', por lo que aparecerá una notificación para contactar con el administrador.
    
    \imagenAncho{Login_forgot_password}{Login - \textit{Forgot password}}{0.5}
    
    \item Si el usuario quiere ver la contraseña que ha introducido puede hacer click en el icono del ojo que se encuentra en el campo.
\end{itemize}

\imagenAncho{Login}{Login}{0.5}

\subsection{Consulta de reservas}
Para consultar las reservas es necesario haber introducido como mínimo los filtros ``\textit{Fecha desde}'' y ``\textit{Centro/Departamento}'', aunque por defecto en el campo ``\textit{Fecha desde}'' ya está seleccionada la fecha actual. 

\imagenAncho{Consulta_reservas_vacio}{Consulta de Reservas}{1}

Una vez se han introducido los filtros que se quieren aplicar hay que hacer click en ``\textbf{\textit{Buscar}}'', o si se quieren limpiar todos los campos se puede clicar en ``\textbf{\textit{Limpiar filtros}}''.

\begin{itemize}
    \item Si no se ha rellenado el campo ``\textit{Centro/Departamento}'' se mostrará una notificación de error:
    
    \imagenAncho{Consulta_centro_dep_obligatorio}{Consulta de Reservas - Centro/Departamento obligatorio}{0.5}
    
    \item Si se introduce una ``\textit{Fecha desde}'' mayor que la ``\textit{Fecha hasta}'' se mostrará una notificación de error:
    
    \imagenAncho{Consulta_fecha_desde_mayor}{Consulta de Reservas - Fecha desde mayor que fecha hasta}{0.5}
    
    \item Si se introduce una ``\textit{Hora desde}'' mayor que la ``\textit{Hora hasta}'' se mostrará una notificación de error:
    
    \imagenAncho{Consulta_hora_desde_mayor}{Consulta de Reservas - Hora desde mayor que hora hasta}{0.5}
    
    \item Si no hay reservas que cumplan con los filtros aplicados, se mostrará una notificación:
    
    \imagenAncho{Consulta_reservas_no_coincidencias}{Consulta de Reservas - No hay reservas que cumplan con los filtros aplicados}{0.5}
\end{itemize}

Si la consulta es correcta, se muestran en una tabla las reservas registradas que cumplan con los filtros aplicados.

\begin{landscape}
    Por ejemplo, en la imagen se han buscado las reservas a partir del 17/07/2020 realizadas por ``\textit{Departamento Prueba}'' :

    \imagenAncho{Consulta_reservas_ejemplo}{Consulta de Reservas - Ejemplo}{1.6}
\end{landscape}

\imagenAncho{Consulta_aulas_ejemplo}{Consulta de Disponibilidad de Aulas - Ejemplo}{1}

\subsection{Consulta de disponibilidad de aulas}
Para consultar la disponibilidad de las aulas es necesario haber introducido como mínimo el filtro ``\textit{Centro/Departamento}''. 

\imagenAncho{Consulta_aulas_vacio}{Consulta de Disponibilidad de Aulas}{1}

Una vez se han introducido los filtros que se quieren aplicar hay que hacer click en ``\textbf{\textit{Buscar}}'', o si se quieren limpiar todos los campos se puede clicar en ``\textbf{\textit{Limpiar filtros}}''.

\begin{itemize}
    \item Si no se ha rellenado el campo ``\textit{Centro/Departamento}'' se mostrará una notificación de error:
    
    \imagenAncho{Consulta_centro_dep_obligatorio}{Consulta de Disponibilidad de Aulas - Centro/Departamento obligatorio}{0.5}
    
    \item Si se introduce una ``\textit{Fecha desde}'' mayor que la ``\textit{Fecha hasta}'' se mostrará una notificación de error:
    
    \imagenAncho{Consulta_fecha_desde_mayor}{Consulta de Disponibilidad de Aulas - Fecha desde mayor que fecha hasta}{0.5}
    
    \item Si se introduce una ``\textit{Hora desde}'' mayor que la ``\textit{Hora hasta}'' se mostrará una notificación de error:
    
    \imagenAncho{Consulta_hora_desde_mayor}{Consulta de Disponibilidad de Aulas - Hora desde mayor que hora hasta}{0.5}
    
    \item Si se quiere filtrar por fecha y hora, los filtros ``\textit{Fecha desde}'', ``\textit{Hora desde}'' y ``\textit{Hora hasta}'' son obligatorios, en caso de no haberlos rellenado se mostrará una notificación de error:
    
    \imagenAncho{Consulta_aulas_fecha_horas_obligatorio}{Consulta de Disponibilidad de Aulas - Fecha y horas obligatorias}{0.5}
    
    \item Si no hay aulas disponibles que cumplan con los filtros aplicados, se mostrará una notificación:
    
    \imagenAncho{Consulta_aulas_no_coincidencias}{Consulta de Disponibilidad de Aulas - No hay aulas que cumplan con los filtros aplicados}{0.5}
\end{itemize}

Si la consulta es correcta, se muestran en una tabla las aulas disponibles que cumplan con los filtros aplicados.

Por ejemplo, en la imagen se han buscado las aulas cuya capacidad sea como mínimo de 25 personas:

\imagenAncho{Consulta_aulas_ejemplo}{Consulta de Disponibilidad de Aulas - Ejemplo}{1}

\subsection{Administrador}
A continuación se detalla el uso de las ventanas cuando se accede como usuario administrador del sistema.

Las ventanas a las que puede que tiene acceso se muestran en el menú lateral, enfatizándose la ventana en la que se encuentra:

\imagenAncho{Admin_menu_lateral}{Menú lateral del usuario administrador}{0.3}

\subsubsection{Histórico de reservas}
Al acceder al histórico de reservas se mostrará en la parte superior un formulario para poder filtrar por fecha y en la parte inferior un listado con todas las operaciones realizadas sobre las reservas, filtradas por defecto a partir de la fecha actual.

Si no hay operaciones que mostrar, aparecerá una notificación:

\imagenAncho{Historico_no_operaciones}{Histórico de Reservas - No hay operaciones}{0.5}

\imagenAncho{Historico}{Histórico de Reservas}{1}

Si se quiere realizar un \textbf{filtrado de las operaciones}, se rellenarán los campos correspondientes y se clicará en el botón ``\textit{Buscar}'', o si se quieren vaciar los campos se clicará en ``\textit{Limpiar filtros}''.

\begin{itemize}
    \item Si se introduce una ``\textit{Fecha desde}'' mayor que la ``\textit{Fecha hasta}'' se mostrará una notificación de error:
    
    \imagenAncho{Consulta_fecha_desde_mayor}{Gestión de Reservas - Fecha desde mayor que fecha hasta}{0.5}
    
    \item Si no hay operaciones realizadas en las fechas seleccionadas se mostrará la notificación.
\end{itemize}

\subsubsection{Mantenimiento de centros y departamentos}
Cuando se accede al mantenimiento de centros y departamentos, si aún no hay ninguno registrado, se muestran simplemente la barra de herramientas (filtro de búsqueda y botones) y la cabecera de la tabla.

\imagenAncho{Mant_propietarios_vacio}{Mantenimiento de Centros y Departamentos vacío}{1}

Para \textbf{añadir un centro} nuevo se hace click en el botón ``\textbf{\textit{Añadir centro}}'':
\begin{itemize}
    \item Se rellenan los campos de la ventana.
    \item Se puede guardar clicando en el botón ``\textbf{\textit{Guardar}}'' o volver a la ventana anterior, sin guardar los cambios, clicando en ``\textbf{\textit{Cancelar}}''.
\end{itemize}

Una vez añadido, se redirige a la ventana de mantenimiento de centros y departamentos, donde se puede ver el nuevo centro en el listado junto a los datos de su responsable:

\imagenAncho{Mant_propietarios_anadir_centro}{Mantenimiento de Centros y Departamentos - Añadir centro}{1}

Para \textbf{añadir un departamento} el proceso es el mismo, solo que haciendo click en ``\textbf{\textit{Añadir departamento}}''.

Si se quieren \textbf{modificar los datos de un centro o departamento, o eliminarlo}, hay que hacer click en la fila correspondiente a éste. Se abrirá el formulario de edición con los valores actuales, manteniendo el campo del ID deshabilitado (no se puede modificar):

\imagenAncho{Mant_propietarios_editar}{Mantenimiento de Centros y Departamentos - Modificar o eliminar un centro o departamento}{1}

Se pueden guardar los cambios clicando en el botón ``\textbf{\textit{Guardar}}'', o volver a la ventana anterior clicando en ``\textbf{\textit{Cancelar}}''. Al guardar, se actualizarán los valores correspondientes en la tabla.

Se puede eliminar el centro o departamento clicando en el botón ``\textbf{\textit{Eliminar}}''. Aparecerá un cuadro para confirmar la acción, clicando en el botón ``\textbf{\textit{Confirmar}}'', o cancelarla clicando en ``\textbf{\textit{Cancelar}}''.
        
\imagenAncho{Mant_propietarios_eliminar}{Mantenimiento de Centros y Departamentos - Confirmación de eliminación}{1}

Se pueden filtrar los datos de la tabla introduciendo el nombre del centro o departamento, o parte de este, en el campo ``\textit{Buscar por nombre\dots}'' de la barra de herramientas.

\imagenAncho{Mant_propietarios_filtrar_1}{Mantenimiento de Centros y Departamentos - Antes de filtrar por nombre}{1}
\imagenAncho{Mant_propietarios_filtrar_2}{Mantenimiento de Centros y Departamentos - Después de filtrar por nombre}{1}

\subsubsection{Mantenimiento de aulas}
Cuando se accede al mantenimiento de aulas, si no se ha seleccionado ninguno o si el centro o departamento no tiene aulas, se muestran simplemente la barra de herramientas (filtro de búsqueda y botones) y la cabecera de la tabla.

\imagenAncho{Mant_aulas_vacio}{Mantenimiento de Aulas vacío}{1}

Para \textbf{añadir un aula} nuevo se hace click en el botón ``\textbf{\textit{Añadir aula}}'':
\begin{itemize}
    \item Se rellenan los campos de la ventana.
    \item Se puede guardar clicando en el botón ``\textbf{\textit{Guardar}}'' o volver a la ventana anterior, sin guardar los cambios, clicando en ``\textbf{\textit{Cancelar}}''.
\end{itemize}

\imagenAncho{Mant_aulas_anadir}{Mantenimiento de Aulas - Añadir aula}{1}

Una vez añadido, se redirige a la ventana de mantenimiento de aulas, donde se puede ver el nuevo aula en el listado seleccionando previamente su propietario en el desplegable:

\imagenAncho{Mant_aulas_propietario}{Mantenimiento de Aulas - Nuevo aula del propietario}{1}

Si se quieren \textbf{modificar los datos de un aula, o eliminarlo}, hay que hacer click en la fila correspondiente a éste. Se abrirá el formulario de edición con los valores actuales, manteniendo el campo del centro en el que se encuentra (\textit{Centro}) deshabilitado (no se puede modificar):

\imagenAncho{Mant_aulas_editar}{Mantenimiento de Aulas - Modificar o eliminar un aula}{1}

Se pueden guardar los cambios clicando en el botón ``\textbf{\textit{Guardar}}'', o volver a la ventana anterior clicando en ``\textbf{\textit{Cancelar}}''. Al guardar, se actualizarán los valores correspondientes en la tabla.

Se puede eliminar el aula clicando en el botón ``\textbf{\textit{Eliminar}}''. Aparecerá un cuadro para confirmar la acción, clicando en el botón ``\textbf{\textit{Confirmar}}'', o cancelarla clicando en ``\textbf{\textit{Cancelar}}''.
        
\imagenAncho{Mant_aulas_eliminar}{Mantenimiento de Aulas - Confirmación de eliminación}{1}

\subsubsection{Mantenimiento de usuarios}
Cuando se accede al mantenimiento de usuarios se muestra la barra de herramientas (filtro de búsqueda y botones) y una tabla con los usuarios registrados.

\imagenAncho{Mant_usuarios_principal}{Mantenimiento de Usuarios}{1}

Para añadir un nuevo usuario se hace click en el botón ``\textbf{\textit{Añadir usuario}}'':
\begin{itemize}
    \item Se rellenan los campos de la ventana.
        \item Si se introduce un correo con formato incorrecto se muestra un mensaje bajo el campo.
        
        \imagenAncho{Mant_usuarios_correo_incorrecto}{Mantenimiento de Usuarios - Correo incorrecto}{0.7}
        
        \item Si se introduce una contraseña con menos de 5 caracteres se muestra un mensaje bajo el campo.
        
        \imagenAncho{Mant_usuarios_contrasena_incorrecta}{Mantenimiento de Usuarios - Contraseña incorrecta}{0.7}
        
        \item Si se introduce un teléfono de menos, o más, de 9 dígitos se muestra un mensaje bajo el campo.
        
        \imagenAncho{Mant_usuarios_telefono_incorrecto}{Mantenimiento de Usuarios - Teléfono incorrecto}{0.7}
        
        \item El \textit{check} ``\textit{Bloqueado}'' es opcional, indica si el usuario se puede eliminar (no marcado) o no (marcado).
    
    \item Se pueden guardar los cambios clicando en el botón ``\textbf{\textit{Guardar}}'', o volver a la ventana anterior clicando en ``\textbf{\textit{Cancelar}}''.
\end{itemize}

\imagenAncho{Mant_usuarios_anadir}{Mantenimiento de Usuarios - Añadir usuario}{1}

Una vez añadido, se redirige a la ventana de mantenimiento de usuarios, donde se puede ver el nuevo usuario en el listado:

\imagenAncho{Mant_usuarios_principal_nuevo}{Mantenimiento de Usuarios - Nuevo usuario}{1}

Para \textbf{bloquear, desbloquear o eliminar un usuario}, hay que hacer click en la fila correspondiente a éste. Se abrirá el formulario con todos los campos deshabilitados (sólo puede modificar sus datos el propio usuario), exceptuando el \textit{check} ``\textit{Bloqueado}''.

\imagenAncho{Mant_usuarios_editar_eliminar}{Mantenimiento de Usuarios - Bloquear, desbloquear o eliminar usuario}{1}

Se puede bloquear/desbloquear el usuario marcando/desmarcando el \textit{check}, aparecerá un cuadro para confirmar la acción, clicando en el botón ``\textbf{\textit{Confirmar}}'', o cancelarla clicando en ``\textbf{\textit{Cancelar}}''. Los cambios se pueden guardar clicando en el botón ``\textbf{\textit{Guardar}}''.

\imagenAncho{Mant_usuarios_bloquear}{Mantenimiento de Usuarios - Confirmación de bloqueo/desbloqueo de usuario}{1}
    
Se puede eliminar el usuario clicando en el botón ``\textbf{\textit{Eliminar}}''.
    \begin{itemize}
        \item Si el usuario está bloqueado aparecerá una notificación de error:
        
        \imagenAncho{Mant_usuarios_eliminar_usuario_bloqueado}{Mantenimiento de Usuarios - Notificación de error al eliminar usuario bloqueado}{1}
        
        \item Si no, aparecerá un cuadro para confirmar la acción, clicando en el botón ``\textbf{\textit{Confirmar}}'', o cancelarla clicando en ``\textbf{\textit{Cancelar}}''.
        
        \imagenAncho{Mant_usuarios_eliminar}{Mantenimiento de Usuarios - Confirmación de eliminación}{1}
    \end{itemize}
    
Se pueden filtrar los datos de la tabla introduciendo el nombre o apellidos del usuario, o parte de estos, en el campo ``\textit{Buscar por nombre o apellidos\dots}'' de la barra de herramientas.

\imagenAncho{Mant_usuarios_filtro_1}{Mantenimiento de Usuarios - Antes de filtrar por nombre}{1}
\imagenAncho{Mant_usuarios_filtro_2}{Mantenimiento de Usuarios - Después de filtrar por nombre}{1}

\subsubsection{Perfil}
El administrador puede modificar los datos de su perfil y clicar en el botón ``\textbf{\textit{Guardar}}'' para guardar los cambios. 

\begin{itemize}
    \item Se realizarán las mismas comprobaciones que en el Mantenimiento de usuarios, exceptuando la contraseña, que debe estar formada por dígitos, minúsculas y mayúsculas y contener como mínimo 6 caracteres.
    
    \imagenAncho{Perfil_contrasena_incorrecta}{Perfil de Usuario - Contraseña incorrecta}{0.7}
    
    \item Si se ha dejado algún campo en blanco se mostrará una notificación de error:
    
    \imagenAncho{Error_campos_obligatorios}{Perfil de Usuario - Notificación de error campos obligatorios}{0.3}
\end{itemize}

\imagenAncho{Perfil_admin}{Perfil de Usuario - Perfil administrador}{1}

\subsection{Responsable}
A continuación se detalla el uso de las ventanas cuando se accede como usuario con rol de responsable.

Las ventanas a las que puede que tiene acceso se muestran en el menú lateral, enfatizándose la ventana en la que se encuentra:

\imagenAncho{Responsable_menu_lateral}{Menú lateral del usuario con rol de responsable}{0.3}

\subsubsection{Reserva de aulas}
Cuando se accede a la reserva de aulas lo único que se ve es el botón ``\textit{Añadir reserva}''.

\imagenAncho{Reserva}{Reserva de Aulas}{0.3}

Cuando se pulsa en el botón para añadir una nueva reserva, se abre el formulario.

\imagenAncho{Reserva_formulario}{Reserva de Aulas - Formulario de reserva}{1}

Si lo que se quiere hacer es una \textbf{reserva por rango de fechas}, se debe marcar el \textit{check} correspondiente que se encuentra al lado de los botones. Una vez marcado se habilitan los campos ``\textit{Fecha fin}'' y ``\textit{Día de la semana}''. ~\footnote{El campo ``\textit{Día de la semana}'' está pensado por si se quieren hacer reservas del tipo ``todos los Lunes'' de \textit{x} fecha a \textit{y} fecha.}

\imagenAncho{Reserva_formulario_rango_fechas}{Reserva de Aulas - Formulario de reserva por rango de fechas}{1}

Una vez se han rellenado todos los campos se hace click en el botón ``\textit{Reservar}'' para registrar la reserva.

\begin{itemize}
    \item Si se ha dejado algún campo en blanco se mostrará una notificación de error:
    
    \imagenAncho{Error_campos_obligatorios}{Reserva de Aulas - Notificación de error campos obligatorios}{0.5}
    
     \item Si se introduce una ``\textit{Fecha inicio}'' mayor que la ``\textit{Fecha fin}'', en caso de reservas por rango de fechas, se mostrará una notificación de error:
    
    \imagenAncho{Reserva_fecha_inicio_mayor}{Reserva de Aulas - Fecha inicio mayor que fecha fin en reservas por rango de fechas}{0.5}
    
    \item Si el aula no se encuentra disponible en la fecha, o fechas, y horas seleccionada, mostrará una notificación de error:
    
    \imagenAncho{Reserva_aula_no_disponible}{Reserva de Aulas - Aula no disponible}{0.5}
\end{itemize}

Una vez registrada, se redirige a la ventana de reserva de aulas.

Un ejemplo de reserva sería el siguiente:

\imagenAncho{Reserva_ejemplo}{Reserva de Aulas - Ejemplo}{1}

\subsubsection{Gestión de reservas}
Al acceder a la gestión de reservas se mostrará en la parte superior un formulario para poder realizar un filtrado y en la parte inferior un listado con todas las reservas registradas a partir de la fecha actual de todos los centros y departamentos bajo responsabilidad del usuario.

Si no hay reservas registradas que mostrar, aparecerá una notificación:

\imagenAncho{Gestion_reservas_no_reservas}{Gestión de Reservas - No hay reservas}{0.5}

Si se quiere realizar un \textbf{filtrado de las reservas mostradas}, se rellenarán los campos correspondientes y se clicará en el botón ``\textit{Buscar}'', o si se quieren vaciar los campos se clicará en ``\textit{Limpiar filtros}''.

Una vez se han introducido los filtros que se quieren aplicar hay que hacer click en ``\textbf{\textit{Buscar}}'', o si se quieren limpiar todos los campos se puede clicar en ``\textbf{\textit{Limpiar filtros}}''.

\begin{itemize}
    \item Si se introduce una ``\textit{Fecha desde}'' mayor que la ``\textit{Fecha hasta}'' se mostrará una notificación de error:
    
    \imagenAncho{Consulta_fecha_desde_mayor}{Gestión de Reservas - Fecha desde mayor que fecha hasta}{0.5}
    
    \item Si se introduce una ``\textit{Hora desde}'' mayor que la ``\textit{Hora hasta}'' se mostrará una notificación de error:
    
    \imagenAncho{Consulta_hora_desde_mayor}{Gestión de Reservas - Hora desde mayor que hora hasta}{0.5}
    
    \item Si no hay reservas que cumplan con los filtros aplicados, se mostrará una notificación:
    
    \imagenAncho{Gestion_reservas_no_reservas}{Gestión de Reservas - No hay reservas (filtrado)}{0.5}
\end{itemize}

Las reservas sólo se pueden modificar de una en una, por lo que si se quiere cambiar algún valor de una reserva por rango de fechas se deberán eliminar todas las reservas correspondientes al rango y volver a registrar una nueva reserva. 

Esta información está disponible en la ventana, clicando en el botón con el icono de información se abrirá un cuadro con la misma y un botón para cerrarle:

\imagenAncho{Gestion_reservas_info}{Gestión de Reservas - Información}{1}

Para \textbf{\textit{modificar una reserva}} se debe marcar el \textit{check} de la fila correspondiente y clicar en el botón ``\textbf{\textit{Modificar}}''.

\begin{itemize}
    \item Si no se ha seleccionado ninguna reserva, se mostrará una notificación de error:
    
    \imagenAncho{Gestion_reservas_no_seleccion}{Gestión de Reservas - No se ha seleccionado ninguna reserva}{0.5}
    
    \item Si se ha seleccionado más de una reserva, se mostrará una notificación de error:
    
    \imagenAncho{Gestion_reservas_editar_unica_reserva}{Gestión de Reservas - Solo se puede modificar una reserva}{0.5}
\end{itemize}

\imagenAncho{Gestion_reservas_editar}{Gestión de Reservas - Modificar}{0.7}

Tras clicar en ``\textit{Modificar}'' se abrirá el formulario de edición con los valores actuales de la reserva. Una vez se hayan modificado los campos deseados se hace click en ``\textbf{\textit{Guardar}}''. Antes de registrar las reservas se realizarán las mismas comprobaciones que en la reserva (en cuanto a fechas, horas, campos obligatorios\dots).

\imagenAncho{Gestion_reservas_editar_2}{Gestión de Reservas - Formulario edición}{1}

Tras modificar los valores de la reserva y guardarlos correctamente, se mostrará en la tabla con los valores actualizados.

\imagenAncho{Gestion_reservas_editar_3}{Gestión de Reservas - Listado con reserva modificada}{1}

Para \textbf{\textit{eliminar una o varias reservas}} se marcan los \textit{checks} correspondientes en la tabla y se hace click en el botón ``\textit{Eliminar}''. 

\begin{itemize}
    \item Al igual que ocurría en la modificación, si no se ha seleccionado ninguna fila se mostrará la notificación de error.
    
    \item Aparecerá un cuadro para confirmar la acción, clicando en el botón  ``\textbf{\textit{Confirmar}}'', o cancelarla clicando en ``\textbf{\textit{Cancelar}}''.
        
    \imagenAncho{Gestion_reservas_eliminar}{Gestión de Reservas - Eliminar reserva/s}{1}
\end{itemize}

Si se confirma el borrado se actualizará el listado de las reservas.

\imagenAncho{Gestion_reservas_eliminar_2}{Gestión de Reservas - Listado de reservas tras el borrado}{1}

\subsubsection{Perfil}
El usuario con rol de responsable puede modificar los datos de su perfil y clicar en el botón ``\textbf{\textit{Guardar}}'' para guardar los cambios. ~\footnote{Se realizarán las mismas comprobaciones que en el perfil del administrador.}

Además, podrá visualizar en una tabla los centros y departamentos de los que es responsable:

\imagenAncho{Perfil_responsable}{Perfil de Usuario - Perfil usuario con rol responsable}{1}


\bibliographystyle{plain}
\bibliography{bibliografiaAnexos}

\end{document}
